\input{preamble}

% OK, start here.
%
\begin{document}

\title{Compléments sur la cohomologie des espaces}

\maketitle

\phantomsection
\label{section-phantom}

\tableofcontents




\section{Introduction}
\label{section-introduction}

\noindent
Nous poursuivons dans ce chapitre l'étude commencée dans
Cohomologie des espaces, section \ref{spaces-cohomology-section-introduction}.
On peut aussi voir ce chapitre comme l'analogue, pour les espaces algébriques,
du chapitre consacré à la cohomologie étale des schémas; voir
Cohomologie étale, section \ref{etale-cohomology-section-introduction}.

\medskip\noindent
En fait, ce chapitre vise principalement à transposer dans le langage des
espaces algébriques des résultats déjà démontrés pour les schémas.
Certains de nos résultats se trouvent dans \cite{Kn}.





\section{Conventions}
\label{section-conventions}

\noindent
Nous supposons une fois pour toutes que tous les schémas appartiennent à
un grand site fppf $\Sch_{fppf}$. Tous les anneaux $A$ considérés
ont en outre la propriété que $\Spec(A)$ est (isomorphe à) un
objet de ce grand site.

\medskip\noindent
Soient $S$ un schéma et $X$ un espace algébrique sur $S$.
Dans ce chapitre et dans les suivants, nous noterons $X \times_S X$
le produit de $X$ par lui-même (dans la catégorie des espaces
algébriques sur $S$), plutôt que $X \times X$.







\section{Transfert des résultats établis pour les schémas}
\label{section-api}

\noindent
Nous expliquons brièvement dans cette section comment des résultats concernant
les schémas entraînent des résultats pour les espaces algébriques
(représentables) et les morphismes (représentables) d'espaces algébriques.
On dispose d'énoncés plus forts pour les modules quasi-cohérents
(car la cohomologie étale d'un module quasi-cohérent sur un schéma
coïncide avec sa cohomologie de Zariski); cette question a déjà été
traitée dans Cohomologie des espaces, section
\ref{spaces-cohomology-section-higher-direct-image}.

\medskip\noindent
Soient $S$ un schéma et $X$ un espace algébrique sur $S$.
Supposons maintenant que $X$ soit représentable par le schéma $X_0$
(notation peu commode mais temporaire; nous disons habituellement simplement
que « $X$ est un schéma »). Dans ce cas, $X$ et $X_0$ ont le même petit
site étale:
$$
X_\etale = (X_0)_\etale
$$
Ceci est signalé dans Propriétés des espaces, section
\ref{spaces-properties-section-etale-site}.
De plus, si $f : X \to Y$ est un morphisme d'espaces algébriques représentables
sur $S$ et si $f_0 : X_0 \to Y_0$ est un morphisme de schémas
représentant $f$, alors les morphismes induits entre les petits topos
étales coïncident:
$$
\xymatrix{
\Sh(X_\etale) \ar[rr]_{f_{small}} \ar@{=}[d] & &
\Sh(Y_\etale) \ar@{=}[d] \\
\Sh((X_0)_\etale) \ar[rr]^{(f_0)_{small}} & &
\Sh((Y_0)_\etale)
}
$$
Voir Propriétés des espaces, lemme
\ref{spaces-properties-lemma-functoriality-etale-site} et
Topologies, lemme \ref{topologies-lemma-morphism-big-small-etale}.

\medskip\noindent
Il n'y a donc absolument aucune différence entre la cohomologie étale
d'un schéma et celle de l'espace algébrique correspondant.
Il en va de même des images directes supérieures par des morphismes de schémas.
En effet, si $f : X \to Y$ est un morphisme d'espaces algébriques sur $S$
qui est représentable (par des schémas), les images directes supérieures
$R^if_*\mathcal{F}$ d'un faisceau $\mathcal{F}$ sur $X_\etale$
peuvent se calculer localement pour la topologie étale sur $Y$ (Cohomologie
sur les sites, lemme \ref{sites-cohomology-lemma-higher-direct-images});
ceci ramène donc souvent les calculs et les démonstrations au cas où
$Y$ et $X$ sont des schémas.

\medskip\noindent
Nous utiliserons ce qui précède sans autre mention dans ce chapitre.
Le même fait vaut pour les autres topologies; nous l'énonçons ici
explicitement sous forme d'un lemme pour la cohomologie.

\begin{lemma}
\label{lemma-compare-cohomology-other-topologies}
Soit $S$ un schéma. Soit $\tau \in \{\etale, fppf, ph\}$ (en ajouter ici).
Le foncteur d'inclusion
$$
(\Sch/S)_\tau \longrightarrow (\textit{Espaces}/S)_\tau
$$
est un foncteur cocontinu spécial
(Sites, définition \ref{sites-definition-special-cocontinuous-functor})
et identifie par conséquent les topos.
\end{lemma}

\begin{proof}
Les conditions de Sites, lemme \ref{sites-lemma-equivalence},
se vérifient immédiatement, puisque notre foncteur est pleinement fidèle
et que tout espace algébrique admet un recouvrement étale par des schémas.
\end{proof}







\section{Changement de base propre}
\label{section-proper-base-change}

\noindent
Le théorème de changement de base propre pour les espaces algébriques se déduit,
au prix d'un peu de travail, du théorème de changement de base propre
pour les schémas et du lemme de Chow.

\begin{lemma}
\label{lemma-surjective-proper}
Soit $S$ un schéma. Soit $f : Y \to X$ un morphisme propre surjectif
d'espaces algébriques sur $S$. Soit $\mathcal{F}$ un faisceau sur $X_\etale$.
Alors $\mathcal{F} \to f_*f^{-1}\mathcal{F}$ est injectif et son
image est l'égalisateur des deux applications
$f_*f^{-1}\mathcal{F} \to g_*g^{-1}\mathcal{F}$, où
$g$ est le morphisme structural $g : Y \times_X Y \to X$.
\end{lemma}

\begin{proof}
Pour tout morphisme surjectif $f : Y \to X$ d'espaces algébriques sur $S$,
l'application $\mathcal{F} \to f_*f^{-1}\mathcal{F}$ est injective.
En effet, si $\overline{x}$ est un point géométrique de $X$, choisissons
un point géométrique $\overline{y}$ de $Y$ situé au-dessus de $\overline{x}$
et considérons
$$
\mathcal{F}_{\overline{x}} \to
(f_*f^{-1}\mathcal{F})_{\overline{x}} \to
(f^{-1}\mathcal{F})_{\overline{y}} = \mathcal{F}_{\overline{x}}
$$
Voir Propriétés des espaces, lemme \ref{spaces-properties-lemma-stalk-pullback},
pour la dernière égalité.

\medskip\noindent
Le second énoncé est local sur $X$ pour la topologie étale; nous pouvons donc
supposer, et nous supposons, que $Y$ est un schéma affine.

\medskip\noindent
Choisissons un morphisme propre surjectif $Z \to Y$, où $Z$ est un schéma; voir
Cohomologie des espaces, lemme \ref{spaces-cohomology-lemma-weak-chow}.
Le résultat pour $Z \to X$ entraîne celui pour $Y \to X$.
Comme $Z \to X$ est un morphisme propre surjectif de schémas,
et par conséquent un recouvrement ph
(Topologies, lemme \ref{topologies-lemma-surjective-proper-ph}),
le résultat pour $Z \to X$ découle de Cohomologie étale, lemme
\ref{etale-cohomology-lemma-describe-pullback-pi-ph}
(il est en fait, en un certain sens, équivalent à ce lemme).
\end{proof}

\begin{lemma}
\label{lemma-h0-proper-over-henselian-pair}
Soit $(A, I)$ un couple hensélien. Soit $X$ un espace algébrique sur $A$
tel que le morphisme structural $f : X \to \Spec(A)$ soit propre.
Soit $i : X_0 \to X$ l'inclusion de $X \times_{\Spec(A)} \Spec(A/I)$.
Pour tout faisceau $\mathcal{F}$ sur $X_\etale$, on a
$\Gamma(X, \mathcal{F}) = \Gamma(X_0, i^{-1}\mathcal{F})$.
\end{lemma}

\begin{proof}
Choisissons un morphisme propre surjectif $Y \to X$, où $Y$ est un schéma; voir
Cohomologie des espaces, lemme \ref{spaces-cohomology-lemma-weak-chow}.
Considérons le diagramme
$$
\xymatrix{
\Gamma(X_0, \mathcal{F}_0) \ar[r] &
\Gamma(Y_0, \mathcal{G}_0) \ar@<1ex>[r] \ar@<-1ex>[r] &
\Gamma((Y \times_X Y)_0, \mathcal{H}_0) \\
\Gamma(X, \mathcal{F}) \ar[r] \ar[u] &
\Gamma(Y, \mathcal{G}) \ar@<1ex>[r] \ar@<-1ex>[r] \ar[u] &
\Gamma(Y \times_X Y, \mathcal{H}) \ar[u]
}
$$
Ici $\mathcal{G}$, resp.\ $\mathcal{H}$, est l'image réciproque de
$\mathcal{F}$ sur $Y$, resp.\ sur $Y \times_X Y$, et l'indice $0$
indique le changement de base à $\Spec(A/I)$. D'après le cas des schémas
(Cohomologie étale, lemme
\ref{etale-cohomology-lemma-h0-proper-over-henselian-pair})
nous voyons que les flèches verticales du milieu et de droite sont bijectives.
Le lemme \ref{lemma-surjective-proper} montre qu'il en va de même de celle de gauche.
\end{proof}

\begin{lemma}
\label{lemma-h0-proper-over-henselian-local}
Soit $A$ un anneau local hensélien. Soit $X$ un espace algébrique
sur $A$ tel que $f : X \to \Spec(A)$ soit un morphisme propre.
Soit $X_0 \subset X$ la fibre de $f$ au-dessus du point fermé.
Pour tout faisceau $\mathcal{F}$ sur $X_\etale$, on a
$\Gamma(X, \mathcal{F}) = \Gamma(X_0, \mathcal{F}|_{X_0})$.
\end{lemma}

\begin{proof}
C'est un cas particulier du lemme \ref{lemma-h0-proper-over-henselian-pair}.
\end{proof}

\begin{lemma}
\label{lemma-proper-base-change-f-star}
Soit $S$ un schéma. Soient $f : X \to Y$ et $g : Y' \to Y$
des morphismes d'espaces algébriques sur $S$. Supposons $f$ propre.
Posons $X' = Y' \times_Y X$, de projections $f' : X' \to Y'$ et $g' : X' \to X$.
Soit $\mathcal{F}$ un faisceau quelconque sur $X_\etale$. Alors
$g^{-1}f_*\mathcal{F} = f'_*(g')^{-1}\mathcal{F}$.
\end{lemma}

\begin{proof}
La question est locale pour la topologie étale sur $Y'$. Choisissons un schéma
$V$ et un morphisme étale surjectif $V \to Y$. Choisissons un schéma $V'$ et un
morphisme étale surjectif $V' \to V \times_Y Y'$. Nous pouvons alors remplacer
$Y'$ par $V'$ et $Y$ par $V$, donc supposer que $Y$ et $Y'$ sont des schémas.
Nous pouvons ensuite travailler localement pour la topologie de Zariski sur
$Y$ et $Y'$, et donc supposer que $Y$ et $Y'$ sont des schémas affines.

\medskip\noindent
Supposons que $Y$ et $Y'$ soient des schémas affines. Choisissons un morphisme
propre surjectif $h_1 : X_1 \to X$, où $X_1$ est un schéma; voir
Cohomologie des espaces, lemme \ref{spaces-cohomology-lemma-weak-chow}.
Posons $X_2 = X_1 \times_X X_1$ et notons
$h_2 : X_2 \to X$ le morphisme structural. Observons qu'il s'agit d'un schéma.
D'après le cas des schémas
(Cohomologie étale, lemme
\ref{etale-cohomology-lemma-proper-base-change-f-star})
nous savons que le lemme est vrai pour les diagrammes cartésiens
$$
\vcenter{
\xymatrix{
X'_1 \ar[r] \ar[d] & X_1 \ar[d] \\
Y' \ar[r] & Y
}
}
\quad\text{et}\quad
\vcenter{
\xymatrix{
X'_2 \ar[r] \ar[d] & X_2 \ar[d] \\
Y' \ar[r] & Y
}
}
$$
et les faisceaux $\mathcal{F}_i = (X_i \to X)^{-1}\mathcal{F}$.
D'après le lemme \ref{lemma-surjective-proper}, nous avons une suite exacte
$0 \to \mathcal{F} \to h_{1, *}\mathcal{F}_1 \to h_{2, *}\mathcal{F}_2$
et de même pour $(g')^{-1}\mathcal{F}$, puisque
$X'_2 = X'_1 \times_{X'} X'_1$. Nous en concluons que le
lemme est vrai (certains détails sont omis).
\end{proof}

\noindent
Soit $S$ un schéma.
Soit $f : Y \to X$ un morphisme d'espaces algébriques sur $S$. Soit
$\overline{x} : \Spec(k) \to S$ un point géométrique. La fibre
de $f$ en $\overline{x}$ est l'espace algébrique
$Y_{\overline{x}} = \Spec(k) \times_{\overline{x}, X} Y$ sur $\Spec(k)$.
Si $\mathcal{F}$ est un faisceau sur $Y_\etale$, notons
$\mathcal{F}_{\overline{x}} = p^{-1}\mathcal{F}$
l'image réciproque de $\mathcal{F}$ sur $(Y_{\overline{x}})_\etale$.
Ici $p : Y_{\overline{x}} \to Y$ est la projection.
Dans la suite, nous considérerons l'ensemble
$\Gamma(Y_{\overline{x}}, \mathcal{F}_{\overline{x}})$.

\begin{lemma}
\label{lemma-proper-pushforward-stalk}
Soit $S$ un schéma.
Soit $f : Y \to X$ un morphisme propre d'espaces algébriques sur $S$. Soit
$\overline{x} \to X$ un point géométrique.
Pour tout faisceau $\mathcal{F}$ sur $Y_\etale$,
l'application canonique
$$
(f_*\mathcal{F})_{\overline{x}} \longrightarrow
\Gamma(Y_{\overline{x}}, \mathcal{F}_{\overline{x}})
$$
est bijective.
\end{lemma}

\begin{proof}
C'est un cas particulier du lemme \ref{lemma-proper-base-change-f-star}.
\end{proof}

\begin{theorem}
\label{theorem-proper-base-change}
Soit $S$ un schéma. Soit
$$
\xymatrix{
X' \ar[r]_{g'} \ar[d]_{f'} & X \ar[d]^f \\
Y' \ar[r]^g & Y
}
$$
un carré cartésien d'espaces algébriques sur $S$.
Supposons $f$ propre.
Soit $\mathcal{F}$ un faisceau abélien de torsion sur $X_\etale$.
Alors le morphisme de changement de base
$$
g^{-1}Rf_*\mathcal{F} \longrightarrow Rf'_*(g')^{-1}\mathcal{F}
$$
est un isomorphisme.
\end{theorem}

\begin{proof}
Cette démonstration reprend certains des arguments de la démonstration du
théorème de changement de base propre pour les schémas. Voir
Cohomologie étale, section \ref{etale-cohomology-section-proper-base-change},
pour davantage de détails.

\medskip\noindent
L'énoncé est local pour la topologie étale sur $Y'$ et sur $Y$; nous pouvons
donc supposer que $Y$ et $Y'$ sont tous deux des schémas affines. Observons que
ceci démontre en particulier le théorème lorsque $f$ est représentable
(nous l'utiliserons plus bas).

\medskip\noindent
Pour tout $n \geq 1$, soit $\mathcal{F}[n]$ le sous-faisceau des sections
de $\mathcal{F}$ annulées par $n$. Alors
$\mathcal{F} = \colim \mathcal{F}[n]$. D'après 
Cohomologie des espaces, lemme \ref{spaces-cohomology-lemma-colimit-cohomology},
les foncteurs $g^{-1}R^pf_*$ et $R^pf'_*(g')^{-1}$ commutent
aux limites inductives filtrantes. Il suffit donc de démontrer le théorème
lorsque $\mathcal{F}$ est annulée par $n$.

\medskip\noindent
Soit $\mathcal{F} \to \mathcal{I}^\bullet$ une résolution par des
faisceaux injectifs de $\mathbf{Z}/n\mathbf{Z}$-modules.
Observons que
$g^{-1}f_*\mathcal{I}^\bullet = f'_*(g')^{-1}\mathcal{I}^\bullet$
d'après le lemme \ref{lemma-proper-base-change-f-star}.
En appliquant le lemme d'acyclicité de Leray
(Catégories dérivées, lemme \ref{derived-lemma-leray-acyclicity}),
nous concluons qu'il suffit de démontrer
$R^pf'_*(g')^{-1}\mathcal{I}^m = 0$ pour $p > 0$ et $m \in \mathbf{Z}$.

\medskip\noindent
Choisissons un morphisme propre surjectif
$h : Z \to X$, où $Z$ est un schéma; voir
Cohomologie des espaces, lemme \ref{spaces-cohomology-lemma-weak-chow}.
Choisissons une application injective $h^{-1}\mathcal{I}^m \to \mathcal{J}$,
où $\mathcal{J}$ est un faisceau injectif de
$\mathbf{Z}/n\mathbf{Z}$-modules sur $Z_\etale$.
Comme $h$ est surjectif, l'application $\mathcal{I}^m \to h_*\mathcal{J}$
est injective (voir le lemme \ref{lemma-surjective-proper}).
Puisque $\mathcal{I}^m$ est injectif, nous voyons que $\mathcal{I}^m$
est un facteur direct de
$h_*\mathcal{J}$. Il suffit donc de démontrer l'annulation voulue
pour $h_*\mathcal{J}$.

\medskip\noindent
Notons $h'$ son changement de base par $g$ et
$g'' : Z' \to Z$ la projection. Il existe une suite spectrale
$$
E_2^{p, q} = R^pf'_* R^qh'_* (g'')^{-1}\mathcal{J}
$$
qui aboutit à $R^{p + q}(f' \circ h')_*(g'')^{-1}\mathcal{J}$.
Comme $h$ et $f \circ h$ sont représentables (par des schémas),
nous savons que le résultat voulu vaut pour eux. La suite spectrale montre donc
que $E_2^{p, q} = 0$ pour $q > 0$ et que
$R^{p + q}(f' \circ h')_*(g'')^{-1}\mathcal{J} = 0$
pour $p + q > 0$. Il s'ensuit que $E_2^{p, 0} = 0$ pour $p > 0$.
Or
$$
E_2^{p, 0} = R^pf'_* h'_* (g'')^{-1}\mathcal{J} =
R^pf'_* (g')^{-1}h_*\mathcal{J}
$$
d'après le lemme \ref{lemma-proper-base-change-f-star}. Ceci achève la démonstration.
\end{proof}

\begin{lemma}
\label{lemma-proper-base-change}
Soit $S$ un schéma. Soit
$$
\xymatrix{
X' \ar[r]_{g'} \ar[d]_{f'} & X \ar[d]^f \\
Y' \ar[r]^g & Y
}
$$
un carré cartésien d'espaces algébriques sur $S$. Supposons $f$ propre.
Soit $E \in D^+(X_\etale)$ à faisceaux de cohomologie de torsion.
Alors le morphisme de changement de base $g^{-1}Rf_*E \to Rf'_*(g')^{-1}E$
est un isomorphisme.
\end{lemma}

\begin{proof}
C'est une conséquence immédiate du théorème de changement de base propre
(théorème \ref{theorem-proper-base-change}), en utilisant les suites
spectrales
$$
E_2^{p, q} = R^pf_*H^q(E)
\quad\text{et}\quad
{E'}_2^{p, q} = R^pf'_*(g')^{-1}H^q(E)
$$
qui aboutissent à $R^nf_*E$ et $R^nf'_*(g')^{-1}E$.
Ces suites spectrales sont construites dans Catégories dérivées, lemme
\ref{derived-lemma-two-ss-complex-functor}.
Certains détails sont omis.
\end{proof}

\begin{lemma}
\label{lemma-proper-base-change-stalk}
Soit $S$ un schéma.
Soit $f : X \to Y$ un morphisme propre d'espaces algébriques.
Soit $\overline{y} \to Y$ un point géométrique.
\begin{enumerate}
\item Pour un faisceau abélien de torsion $\mathcal{F}$ sur $X_\etale$, on a
$(R^nf_*\mathcal{F})_{\overline{y}} =
H^n_\etale(X_{\overline{y}}, \mathcal{F}_{\overline{y}})$.
\item Pour $E \in D^+(X_\etale)$ à faisceaux de cohomologie de torsion, on a
$(R^nf_*E)_{\overline{y}} = H^n_\etale(X_{\overline{y}}, E_{\overline{y}})$.
\end{enumerate}
\end{lemma}

\begin{proof}
Dans l'énoncé, $\mathcal{F}_{\overline{y}}$ désigne l'image réciproque
de $\mathcal{F}$ sur $X_{\overline{y}} = \overline{y} \times_Y X$.
Comme l'image réciproque par $\overline{y} \to Y$ donne la fibre de
$\mathcal{F}$, la première assertion du lemme est un cas particulier du
théorème \ref{theorem-proper-base-change}. La seconde est un cas particulier
du lemme \ref{lemma-proper-base-change}.
\end{proof}

\begin{lemma}
\label{lemma-base-change-separably-closed}
Soit $k'/k$ une extension de corps séparablement clos.
Soit $X$ un espace algébrique propre sur $k$.
Soit $\mathcal{F}$ un faisceau abélien de torsion sur $X$.
Alors l'application $H^q_\etale(X, \mathcal{F}) \to
H^q_\etale(X_{k'}, \mathcal{F}|_{X_{k'}})$ est un isomorphisme
pour $q \geq 0$.
\end{lemma}

\begin{proof}
C'est un cas particulier du théorème \ref{theorem-proper-base-change}.
\end{proof}










\section{Comparaison des grands et petits topos}
\label{section-compare}

\noindent
Soient $S$ un schéma et $X$ un espace algébrique sur $S$.
Dans Topologies sur les espaces, lemme
\ref{spaces-topologies-lemma-at-the-bottom-etale}
nous avons introduit les morphismes de comparaison
$\pi_X : (\textit{Espaces}/X)_\etale \to X_{spaces, \etale}$ et
$i_X : \Sh(X_\etale) \to \Sh((\textit{Espaces}/X)_\etale)$
avec $\pi_X \circ i_X = \text{id}$ comme morphismes de topos et
$\pi_{X, *} = i_X^{-1}$.
Plus généralement, si $f : Y \to X$ est un objet de $(\textit{Espaces}/X)_\etale$,
il existe un morphisme
$i_f : \Sh(Y_\etale) \to \Sh((\textit{Espaces}/X)_\etale)$
tel que $f_{small} = \pi_X \circ i_f$; voir
Topologies sur les espaces, lemmes \ref{spaces-topologies-lemma-put-in-T-etale} et
\ref{spaces-topologies-lemma-morphism-big-small-etale}. Dans
Topologies sur les espaces, remarque
\ref{spaces-topologies-remark-change-topologies-ringed}
nous avons étendu ceux-ci en un morphisme de sites annelés
$$
\pi_X :
((\textit{Espaces}/X)_\etale, \mathcal{O})
\to
(X_{spaces, \etale}, \mathcal{O}_X)
$$
et en des morphismes de topos annelés
$$
i_X :
(\Sh(X_\etale), \mathcal{O}_X)
\to
(\Sh((\textit{Espaces}/X)_\etale), \mathcal{O})
$$
et
$$
i_f :
(\Sh(Y_\etale), \mathcal{O}_Y)
\to
(\Sh((\textit{Espaces}/X)_\etale, \mathcal{O}))
$$
Notons que la restriction $i_X^{-1} = \pi_{X, *}$ (voir
Topologies, définition \ref{topologies-definition-restriction-small-etale})
transforme $\mathcal{O}$ en $\mathcal{O}_X$.
De même, $i_f^{-1}$ transforme $\mathcal{O}$ en $\mathcal{O}_Y$.
Voir Topologies sur les espaces, remarque
\ref{spaces-topologies-remark-change-topologies-ringed}.
Ainsi $i_X^*\mathcal{F} = i_X^{-1}\mathcal{F}$ et
$i_f^*\mathcal{F} = i_f^{-1}\mathcal{F}$ pour tout $\mathcal{O}$-module
$\mathcal{F}$ sur $(\textit{Espaces}/X)_\etale$. En particulier, $i_X^*$ et
$i_f^*$ sont des foncteurs exacts. Le foncteur $i_X^*$ est souvent noté
$\mathcal{F} \mapsto \mathcal{F}|_{X_\etale}$ (ce qui n'entre pas
en conflit avec la notation de Topologies sur les espaces, définition
\ref{spaces-topologies-definition-restriction-small-etale}).

\begin{lemma}
\label{lemma-describe-pullback}
Soit $S$ un schéma. Soit $X$ un espace algébrique sur $S$.
Soit $\mathcal{F}$ un faisceau sur $X_\etale$. Alors
$\pi_X^{-1}\mathcal{F}$ est donné par la règle
$$
(\pi_X^{-1}\mathcal{F})(Y) = \Gamma(Y_\etale, f_{small}^{-1}\mathcal{F})
$$
pour $f : Y \to X$ dans $(\textit{Espaces}/X)_\etale$.
De plus, $\pi_Y^{-1}\mathcal{F}$ satisfait à la condition de faisceau pour les
recouvrements lisses, syntomiques, fppf, fpqc et ph.
\end{lemma}

\begin{proof}
Comme l'image réciproque est transitive et que $f_{small} = \pi_X \circ i_f$
(voir ci-dessus), nous avons
$i_f^{-1} \pi_X^{-1}\mathcal{F} = f_{small}^{-1}\mathcal{F}$.
Ceci donne la description de $\pi_X^{-1}$ annoncée dans le lemme.

\medskip\noindent
Pour démontrer que $\pi_X^{-1}\mathcal{F}$ est un faisceau pour la topologie ph,
il suffit, d'après Topologies sur les espaces, lemme
\ref{spaces-topologies-lemma-characterize-sheaf}
de montrer que, pour tout morphisme propre surjectif
$V \to U$ d'espaces algébriques sur $X$, l'ensemble
$(\pi_X^{-1}\mathcal{F})(U)$ est l'égalisateur des deux applications
$(\pi_X^{-1}\mathcal{F})(V) \to (\pi_X^{-1}\mathcal{F})(V \times_U V)$.
C'est précisément le lemme \ref{lemma-surjective-proper}.

\medskip\noindent
Le cas des recouvrements lisses, syntomiques et fppf se déduit de celui
des recouvrements ph par Topologies sur les espaces, lemme
\ref{spaces-topologies-lemma-zariski-etale-smooth-syntomic-fppf-ph}.

\medskip\noindent
Soit $\mathcal{U} = \{U_i \to U\}_{i \in I}$ un recouvrement fpqc d'espaces
algébriques sur $X$. Soient $s_i \in (\pi_X^{-1}\mathcal{F})(U_i)$ des sections
qui coïncident sur $U_i \times_U U_j$. Il faut démontrer qu'il existe un unique
$s \in (\pi_X^{-1}\mathcal{F})(U)$ qui se restreint en $s_i$ sur $U_i$.
Cas I: $U$ et les $U_i$ sont des schémas. Ce cas découle de
Cohomologie étale, lemme \ref{etale-cohomology-lemma-describe-pullback}.
Cas II: $U$ est un schéma. Choisissons des morphismes étales surjectifs
$T_i \to U_i$, où les $T_i$ sont des schémas. Alors
$\mathcal{T} = \{T_i \to U\}$ est un recouvrement fpqc par des schémas et,
d'après le cas I, le résultat vaut pour $\mathcal{T}$. Nous omettons de vérifier
que ceci entraîne le résultat pour $\mathcal{U}$.
Cas III: cas général. Soit $W \to U$ un morphisme étale surjectif, où $W$ est
un schéma. Alors $\mathcal{W} = \{U_i \times_U W \to W\}$ est un recouvrement
fpqc (par des espaces algébriques) du schéma $W$.
D'après le cas II, le résultat vaut pour $\mathcal{W}$. Nous omettons de vérifier
que ceci entraîne le résultat pour $\mathcal{U}$.
\end{proof}

\begin{lemma}
\label{lemma-compare-injectives}
Soit $S$ un schéma.
Soit $Y \to X$ un morphisme de $(\textit{Espaces}/S)_\etale$.
\begin{enumerate}
\item Si $\mathcal{I}$ est injectif dans
$\textit{Ab}((\textit{Espaces}/X)_\etale)$, alors
\begin{enumerate}
\item $i_f^{-1}\mathcal{I}$ est injectif dans $\textit{Ab}(Y_\etale)$,
\item $\mathcal{I}|_{X_\etale}$ est injectif dans $\textit{Ab}(X_\etale)$,
\end{enumerate}
\item Si $\mathcal{I}^\bullet$ est un complexe K-injectif
dans $\textit{Ab}((\textit{Espaces}/X)_\etale)$, alors
\begin{enumerate}
\item $i_f^{-1}\mathcal{I}^\bullet$ est un complexe K-injectif dans
$\textit{Ab}(Y_\etale)$,
\item $\mathcal{I}^\bullet|_{X_\etale}$ est un complexe K-injectif dans
$\textit{Ab}(X_\etale)$,
\end{enumerate}
\end{enumerate}
Les énoncés correspondants ne valent pas pour les modules.
\end{lemma}

\begin{proof}
Les assertions (1)(b) et (2)(b) découlent formellement du fait que le foncteur
de restriction $\pi_{X, *} = i_X^{-1}$ est adjoint à droite du foncteur exact
$\pi_X^{-1}$; voir Homologie, lemme
\ref{homology-lemma-adjoint-preserve-injectives}, et Catégories dérivées,
lemme \ref{derived-lemma-adjoint-preserve-K-injectives}.

\medskip\noindent
Les assertions (1)(a) et (2)(a) peuvent se voir de deux façons.
Première démonstration: on utilise que $i_f^{-1}$ est adjoint à droite du
foncteur exact $i_{f, !}$. Ce foncteur est construit, pour les faisceaux
d'ensembles, dans Topologies, lemme \ref{topologies-lemma-put-in-T-etale},
et, pour les faisceaux abéliens, dans Modules sur les sites, lemme
\ref{sites-modules-lemma-g-shriek-adjoint}. Il est démontré dans Modules sur
les sites, lemme \ref{sites-modules-lemma-exactness-lower-shriek}, qu'il est exact.
Seconde démonstration: on utilise l'égalité $i_f = i_Y \circ f_{big}$,
établie dans Topologies, lemme
\ref{topologies-lemma-morphism-big-small-etale}. Comme $f_{big}$ est une
localisation, l'image réciproque par ce morphisme préserve les injectifs et les
K-injectifs; voir Cohomologie sur les sites, lemmes
\ref{sites-cohomology-lemma-cohomology-of-open} et
\ref{sites-cohomology-lemma-restrict-K-injective-to-open}.
On conclut alors en appliquant au foncteur $i_Y^{-1}$ les assertions (1)(b)
et (2)(b), déjà démontrées.

\medskip\noindent
Pour un contre-exemple dans le cas des modules, nous renvoyons à
Cohomologie étale, lemme \ref{etale-cohomology-lemma-compare-injectives}.
\end{proof}

\noindent
Soit $S$ un schéma.
Soit $f : Y \to X$ un morphisme d'espaces algébriques sur $S$.
Le diagramme commutatif de Topologies sur les espaces, lemme
\ref{spaces-topologies-lemma-morphism-big-small-etale} (3)
donne un diagramme commutatif de sites annelés
$$
\xymatrix{
(Y_{spaces, \etale}, \mathcal{O}_Y) \ar[d]_{f_{spaces, \etale}} &
((\textit{Espaces}/Y)_\etale, \mathcal{O}) \ar[d]^{f_{big}} \ar[l]^{\pi_Y} \\
(X_{spaces, \etale}, \mathcal{O}_X) &
((\textit{Espaces}/X)_\etale, \mathcal{O}) \ar[l]_{\pi_X}
}
$$
comme on le voit aisément en explicitant les définitions de
$f_{small}^\sharp$, $f_{big}^\sharp$, $\pi_X^\sharp$ et $\pi_Y^\sharp$.
En particulier, cela signifie que
\begin{equation}
\label{equation-compare-big-small}
(f_{big, *}\mathcal{F})|_{X_\etale} =
f_{small, *}(\mathcal{F}|_{Y_\etale})
\end{equation}
pour tout faisceau $\mathcal{F}$ sur $(\textit{Espaces}/Y)_\etale$; si
$\mathcal{F}$ est un faisceau de $\mathcal{O}$-modules, alors
(\ref{equation-compare-big-small})
est un isomorphisme de $\mathcal{O}_X$-modules sur $X_\etale$.

\begin{lemma}
\label{lemma-compare-higher-direct-image}
Soit $S$ un schéma.
Soit $f : Y \to X$ un morphisme d'espaces algébriques sur $S$.
\begin{enumerate}
\item Pour $K$ dans $D((\textit{Espaces}/Y)_\etale)$, on a
$
(Rf_{big, *}K)|_{X_\etale} = Rf_{small, *}(K|_{Y_\etale})
$
dans $D(X_\etale)$.
\item Pour $K$ dans $D((\textit{Espaces}/Y)_\etale, \mathcal{O})$, on a
$
(Rf_{big, *}K)|_{X_\etale} = Rf_{small, *}(K|_{Y_\etale})
$
dans $D(\textit{Mod}(X_\etale, \mathcal{O}_X))$.
\end{enumerate}
Plus généralement, soit $g : X' \to X$ un objet de
$(\textit{Espaces}/X)_\etale$. Considérons le produit fibré
$$
\xymatrix{
Y' \ar[r]_{g'} \ar[d]_{f'} & Y \ar[d]^f \\
X' \ar[r]^g & X
}
$$
Alors
\begin{enumerate}
\item[(3)] Pour $K$ dans $D((\textit{Espaces}/Y)_\etale)$, on a
$i_g^{-1}(Rf_{big, *}K) = Rf'_{small, *}(i_{g'}^{-1}K)$
dans $D(X'_\etale)$.
\item[(4)] Pour $K$ dans $D((\textit{Espaces}/Y)_\etale, \mathcal{O})$, on a
$i_g^*(Rf_{big, *}K) = Rf'_{small, *}(i_{g'}^*K)$
dans $D(\textit{Mod}(X'_\etale, \mathcal{O}_{X'}))$.
\item[(5)] Pour $K$ dans $D((\textit{Espaces}/Y)_\etale)$, on a
$g_{big}^{-1}(Rf_{big, *}K) = Rf'_{big, *}((g'_{big})^{-1}K)$
dans $D((\textit{Espaces}/X')_\etale)$.
\item[(6)] Pour $K$ dans $D((\textit{Espaces}/Y)_\etale, \mathcal{O})$, on a
$g_{big}^*(Rf_{big, *}K) = Rf'_{big, *}((g'_{big})^*K)$
dans $D(\textit{Mod}(X'_\etale, \mathcal{O}_{X'}))$.
\end{enumerate}
\end{lemma}

\begin{proof}
L'assertion (1) découle du lemme \ref{lemma-compare-injectives}
et de (\ref{equation-compare-big-small}),
en choisissant un complexe K-injectif de faisceaux abéliens représentant $K$.

\medskip\noindent
L'assertion (3) découle du lemme \ref{lemma-compare-injectives}
et de Topologies, lemme
\ref{topologies-lemma-morphism-big-small-cartesian-diagram-etale}
en choisissant un complexe K-injectif de faisceaux abéliens représentant $K$.

\medskip\noindent
L'assertion (5) est Cohomologie sur les sites, lemme
\ref{sites-cohomology-lemma-localize-cartesian-square}.

\medskip\noindent
L'assertion (6) est Cohomologie sur les sites, lemme
\ref{sites-cohomology-lemma-localize-cartesian-square-modules}.

\medskip\noindent
L'assertion (2) se démontre comme suit. Nous avons vu ci-dessus que
$\pi_X \circ f_{big} = f_{small} \circ \pi_Y$ comme morphismes de sites
annelés. Nous obtenons donc
$R\pi_{X, *} \circ Rf_{big, *} = Rf_{small, *} \circ R\pi_{Y, *}$
d'après Cohomologie sur les sites, lemme
\ref{sites-cohomology-lemma-derived-pushforward-composition}.
Comme les foncteurs de restriction $\pi_{X, *}$ et $\pi_{Y, *}$
sont exacts, on conclut.

\medskip\noindent
L'assertion (4) découle des assertions (6) et (2), cette dernière étant
appliquée à $f' : Y' \to X'$.
\end{proof}

\noindent
Soit $S$ un schéma. Soit $X$ un espace algébrique sur $S$.
Soit $\mathcal{H}$ un faisceau abélien sur
$(\textit{Espaces}/X)_\etale$. Rappelons que $H^n_\etale(U, \mathcal{H})$
désigne la cohomologie de $\mathcal{H}$ sur un objet
$U$ de $(\textit{Espaces}/X)_\etale$.

\begin{lemma}
\label{lemma-compare-cohomology}
Soit $S$ un schéma.
Soit $f : Y \to X$ un morphisme d'espaces algébriques sur $S$. Alors
\begin{enumerate}
\item Pour $K$ dans $D(X_\etale)$, on a
$H^n_\etale(X, \pi_X^{-1}K) = H^n(X_\etale, K)$.
\item Pour $K$ dans $D(X_\etale, \mathcal{O}_X)$, on a
$H^n_\etale(X, L\pi_X^*K) = H^n(X_\etale, K)$.
\item Pour $K$ dans $D(X_\etale)$, on a
$H^n_\etale(Y, \pi_X^{-1}K) = H^n(Y_\etale, f_{small}^{-1}K)$.
\item Pour $K$ dans $D(X_\etale, \mathcal{O}_X)$, on a
$H^n_\etale(Y, L\pi_X^*K) = H^n(Y_\etale, Lf_{small}^*K)$.
\item Pour $M$ dans $D((\textit{Espaces}/X)_\etale)$, on a
$H^n_\etale(Y, M) = H^n(Y_\etale, i_f^{-1}M)$.
\item Pour $M$ dans $D((\textit{Espaces}/X)_\etale, \mathcal{O})$, on a
$H^n_\etale(Y, M) = H^n(Y_\etale, i_f^*M)$.
\end{enumerate}
\end{lemma}

\begin{proof}
Pour démontrer (5), représentons $M$ par un complexe K-injectif de faisceaux
abéliens, appliquons le lemme \ref{lemma-compare-injectives} et explicitons
les définitions. L'assertion (3) en découle puisque
$i_f^{-1}\pi_X^{-1} = f_{small}^{-1}$. L'assertion (1) est un cas
particulier de (3).

\medskip\noindent
L'assertion (6) découle du très général lemme
\ref{sites-cohomology-lemma-pullback-same-cohomology} de Cohomologie sur les
sites. Ensuite, l'assertion (4) découle de l'égalité
$Lf_{small}^* = i_f^* \circ L\pi_X^*$. L'assertion (2) est un cas particulier
de (4).
\end{proof}

\begin{lemma}
\label{lemma-cohomological-descent-etale}
Soit $S$ un schéma. Soit $X$ un espace algébrique sur $S$.
Pour $K \in D(X_\etale)$, l'application
$$
K \longrightarrow R\pi_{X, *}\pi_X^{-1}K
$$
est un isomorphisme, où
$\pi_X : \Sh((\textit{Espaces}/X)_\etale) \to \Sh(X_\etale)$ est défini ci-dessus.
\end{lemma}

\begin{proof}
Ceci est vrai parce que $\pi_X^{-1}$ et $\pi_{X, *} = i_X^{-1}$ sont tous deux
des foncteurs exacts et que le composé $\pi_{X, *} \circ \pi_X^{-1}$
est le foncteur identité.
\end{proof}

\begin{lemma}
\label{lemma-compare-higher-direct-image-proper}
Soit $S$ un schéma.
Soit $f : Y \to X$ un morphisme propre d'espaces algébriques sur $S$.
Alors on a
\begin{enumerate}
\item $\pi_X^{-1} \circ f_{small, *} = f_{big, *} \circ \pi_Y^{-1}$
comme foncteurs $\Sh(Y_\etale) \to \Sh((\textit{Espaces}/X)_\etale)$,
\item $\pi_X^{-1}Rf_{small, *}K = Rf_{big, *}\pi_Y^{-1}K$
pour $K$ dans $D^+(Y_\etale)$ dont les faisceaux de cohomologie sont de torsion, et
\item $\pi_X^{-1}Rf_{small, *}K = Rf_{big, *}\pi_Y^{-1}K$
pour tout $K$ dans $D(Y_\etale)$ si $f$ est fini.
\end{enumerate}
\end{lemma}

\begin{proof}
Démonstration de (1). Soit $\mathcal{F}$ un faisceau sur $Y_\etale$.
Soit $g : X' \to X$ un objet de $(\textit{Espaces}/X)_\etale$.
Considérons le produit fibré
$$
\xymatrix{
Y' \ar[r]_{f'} \ar[d]_{g'} & X' \ar[d]^g \\
Y \ar[r]^f & X
}
$$
On a alors
$$
(f_{big, *}\pi_Y^{-1}\mathcal{F})(X') =
(\pi_Y^{-1}\mathcal{F})(Y') =
((g'_{small})^{-1}\mathcal{F})(Y')  =
(f'_{small, *}(g'_{small})^{-1}\mathcal{F})(X')
$$
la seconde égalité résultant du lemme \ref{lemma-describe-pullback}.
D'autre part,
$$
(\pi_X^{-1}f_{small, *}\mathcal{F})(X') =
(g_{small}^{-1}f_{small, *}\mathcal{F})(X')
$$
de nouveau d'après le lemme \ref{lemma-describe-pullback}.
Le changement de base propre pour les faisceaux d'ensembles
(lemme \ref{lemma-proper-base-change-f-star}) montre donc que les deux
ensembles sont canoniquement isomorphes. Cet isomorphisme est compatible
avec les applications de restriction et définit un isomorphisme
$\pi_X^{-1}f_{small, *}\mathcal{F} = f_{big, *}\pi_Y^{-1}\mathcal{F}$.
On obtient ainsi un isomorphisme de foncteurs
$\pi_X^{-1} \circ f_{small, *} = f_{big, *} \circ \pi_Y^{-1}$.

\medskip\noindent
Démonstration de (2). Il existe un morphisme canonique de changement de base
$\pi_X^{-1}Rf_{small, *}K \to Rf_{big, *}\pi_Y^{-1}K$
pour tout $K$ dans $D(Y_\etale)$; voir Cohomologie sur les sites, remarque
\ref{sites-cohomology-remark-base-change}.
Pour démontrer que c'est un isomorphisme, il suffit de montrer que l'image
réciproque du morphisme de changement de base par
$i_g : \Sh(X'_\etale) \to \Sh((\Sch/X)_\etale)$ est un isomorphisme pour tout
objet $g : X' \to X$ de $(\Sch/X)_\etale$.
Soient $T', g', f'$ comme au paragraphe précédent.
L'image réciproque du morphisme de changement de base est
\begin{align*}
g_{small}^{-1}Rf_{small, *}K
& =
i_g^{-1}\pi_X^{-1}Rf_{small, *}K \\
& \to
i_g^{-1}Rf_{big, *}\pi_Y^{-1}K \\
& =
Rf'_{small, *}(i_{g'}^{-1}\pi_Y^{-1}K) \\
& =
Rf'_{small, *}((g'_{small})^{-1}K)
\end{align*}
où nous avons utilisé $\pi_X \circ i_g = g_{small}$,
$\pi_Y \circ i_{g'} = g'_{small}$ et le lemme
\ref{lemma-compare-higher-direct-image}.
Cette application est un isomorphisme d'après le théorème de changement de base
propre (lemme \ref{lemma-proper-base-change}), pourvu que $K$ soit borné
inférieurement et que les faisceaux de cohomologie de $K$ soient de torsion.

\medskip\noindent
Démonstration de (3). Si $f$ est fini, les foncteurs
$f_{small, *}$ et $f_{big, *}$ sont exacts. Pour $f_{small}$, ceci découle de
Cohomologie des espaces, lemme
\ref{spaces-cohomology-lemma-finite-higher-direct-image-zero}.
Comme tout changement de base $f'$ de $f$ est encore fini, le lemme
\ref{lemma-compare-higher-direct-image}, assertion (3), montre que
$f_{big, *}$ est lui aussi exact (puisque les foncteurs dérivés supérieurs
sont nuls). Ce cas découle donc de l'assertion (1).
\end{proof}







\section{Comparaison des topologies fppf et étale}
\label{section-fppf-etale}

\noindent
Cette section est l'analogue de Cohomologie étale, section
\ref{etale-cohomology-section-fppf-etale}.

\medskip\noindent
Soient $S$ un schéma et $X$ un espace algébrique sur $S$.
Sur la catégorie $\textit{Espaces}/X$, considérons les topologies fppf
et étale. Le foncteur identité
$(\textit{Espaces}/X)_\etale \to (\textit{Espaces}/X)_{fppf}$
est continu et définit un morphisme de sites
$$
\epsilon_X :
(\textit{Espaces}/X)_{fppf} \longrightarrow (\textit{Espaces}/X)_\etale
$$
par application de Sites, proposition \ref{sites-proposition-get-morphism}.
Notons que $\epsilon_{X, *}$ est le foncteur identité sur les préfaisceaux
sous-jacents et que $\epsilon_X^{-1}$ associe à un faisceau étale son
faisceau associé pour la topologie fppf.
Considérons le morphisme de sites
$$
\pi_X : (\textit{Espaces}/X)_\etale \longrightarrow X_{spaces, \etale}
$$
qui compare les grands et petits sites étales; voir la section \ref{section-compare}.
Le composé détermine un morphisme de sites
$$
a_X = \pi_X \circ \epsilon_X :
(\textit{Espaces}/X)_{fppf}
\longrightarrow
X_{spaces, \etale}
$$
Si $\mathcal{H}$ est un faisceau abélien sur $(\textit{Espaces}/X)_{fppf}$,
nous noterons $H^n_{fppf}(U, \mathcal{H})$ la cohomologie de
$\mathcal{H}$ sur un objet $U$ de $(\textit{Espaces}/X)_{fppf}$.

\begin{lemma}
\label{lemma-comparison-fppf-etale}
Soit $S$ un schéma. Soit $X$ un espace algébrique sur $S$.
\begin{enumerate}
\item Pour $\mathcal{F} \in \Sh(X_\etale)$, on a
$\epsilon_{X, *}a_X^{-1}\mathcal{F} = \pi_X^{-1}\mathcal{F}$
et $a_{X, *}a_X^{-1}\mathcal{F} = \mathcal{F}$.
\item Pour $\mathcal{F} \in \textit{Ab}(X_\etale)$, on a
$R^i\epsilon_{X, *}(a_X^{-1}\mathcal{F}) = 0$ pour $i > 0$.
\end{enumerate}
\end{lemma}

\begin{proof}
On a $a_X^{-1}\mathcal{F} = \epsilon_X^{-1} \pi_X^{-1}\mathcal{F}$.
D'après le lemme \ref{lemma-describe-pullback}, le faisceau étale
$\pi_X^{-1}\mathcal{F}$ est un faisceau pour la topologie fppf; il est donc
égal à $a_X^{-1}\mathcal{F}$ (puisque l'image réciproque par $\epsilon_X$
est donnée par le faisceau associé pour la topologie fppf).
Rappelons en outre que $\epsilon_{X, *}$ est le foncteur identité
sur les préfaisceaux sous-jacents.
L'assertion (1) résulte alors immédiatement de la description explicite de
$\pi_X^{-1}$ donnée au lemme \ref{lemma-describe-pullback}.

\medskip\noindent
Nous démontrerons (2) en nous ramenant au cas des schémas; voir l'assertion (1)
de Cohomologie étale, lemme
\ref{etale-cohomology-lemma-V-C-all-n-etale-fppf}.
Cette méthode « fonctionne clairement », puisque tout espace algébrique est
localement un schéma pour la topologie étale. Les détails sont donnés ci-dessous,
mais nous invitons vivement le lecteur à omettre la démonstration.

\medskip\noindent
Pour un faisceau abélien $\mathcal{H}$ sur $(\textit{Espaces}/X)_{fppf}$,
l'image directe supérieure $R^p\epsilon_{X, *}\mathcal{H}$ est le faisceau
associé au préfaisceau $U \mapsto H^p_{fppf}(U, \mathcal{H})$
sur $(\textit{Espaces}/X)_\etale$. Voir Cohomologie sur les sites, lemme
\ref{sites-cohomology-lemma-higher-direct-images}.
Comme tout objet de $(\textit{Espaces}/X)_\etale$ admet un recouvrement par des
schémas, il suffit de démontrer que, pour tout schéma $U/X$ et tout
$\xi \in H^p_{fppf}(U, a_X^{-1}\mathcal{F})$, il existe un recouvrement étale
$\{U_i \to U\}$ tel que la restriction de $\xi$ à chaque $U_i$ soit nulle.
On a
\begin{align*}
H^p_{fppf}(U, a_X^{-1}\mathcal{F})
& =
H^p((\textit{Espaces}/U)_{fppf}, (a_X^{-1}\mathcal{F})|_{\textit{Espaces}/U}) \\
& =
H^p((\Sch/U)_{fppf}, (a_X^{-1}\mathcal{F})|_{\Sch/U})
\end{align*}
où la seconde identification est le lemme
\ref{lemma-compare-cohomology-other-topologies}, et la première un fait général
sur la restriction (Cohomologie sur les sites, lemme
\ref{sites-cohomology-lemma-cohomology-of-open}).
En comparant le premier paragraphe et le résultat correspondant dans le cas
des schémas (Cohomologie étale, lemme
\ref{etale-cohomology-lemma-describe-pullback-pi-fppf})
nous concluons que le faisceau $(a_X^{-1}\mathcal{F})|_{\Sch/U}$
coïncide avec l'image réciproque par la « version pour les schémas de $a_U$ ».
Nous pouvons donc trouver un recouvrement étale $\{U_i \to U\}$ tel que notre
classe s'annule dans
$H^p((\Sch/U_i)_{fppf}, (a_X^{-1}\mathcal{F})|_{\Sch/U_i})$
pour tout $i$; voir Cohomologie étale, lemme
\ref{etale-cohomology-lemma-V-C-all-n-etale-fppf}
(l'énoncé précis à utiliser ici est que $V_n$ vaut pour tout $n$, ce qui est
l'assertion (2) dans le cas des schémas).
En revenant en arrière (au moyen des mêmes formules que ci-dessus, appliquées
cette fois à $U_i$), nous concluons que la restriction de $\xi$ à $U_i$
est nulle, comme souhaité.
\end{proof}

\noindent
Le travail difficile effectué dans le cas des schémas montre maintenant que
les cohomologies étale et fppf coïncident pour les faisceaux provenant du
petit site étale.

\begin{lemma}
\label{lemma-cohomological-descent-etale-fppf}
Soit $S$ un schéma. Soit $X$ un espace algébrique sur $S$.
Pour $K \in D^+(X_\etale)$, les applications
$$
\pi_X^{-1}K \longrightarrow R\epsilon_{X, *}a_X^{-1}K
\quad\text{et}\quad
K \longrightarrow Ra_{X, *}a_X^{-1}K
$$
sont des isomorphismes, où
$a_X : \Sh((\textit{Espaces}/X)_{fppf}) \to \Sh(X_\etale)$ est défini ci-dessus.
\end{lemma}

\begin{proof}
Nous démontrons seulement le second énoncé; le premier est plus facile et se
démontre exactement de la même manière. On se ramène immédiatement au cas où
$K$ est donné par un seul faisceau abélien. En effet, représentons $K$ par un
complexe borné inférieurement $\mathcal{F}^\bullet$. Le cas d'un faisceau donne
$\mathcal{F}^n = a_{X, *} a_X^{-1} \mathcal{F}^n$
et les faisceaux $R^qa_{X, *}a_X^{-1}\mathcal{F}^n$ sont nuls pour $q > 0$.
Le lemme d'acyclicité de Leray (Catégories dérivées, lemme
\ref{derived-lemma-leray-acyclicity}), appliqué à
$a_X^{-1}\mathcal{F}^\bullet$ et au foncteur $a_{X, *}$, permet de conclure.
Supposons désormais $K = \mathcal{F}$.

\medskip\noindent
Le lemme \ref{lemma-comparison-fppf-etale} donne
$a_{X, *}a_X^{-1}\mathcal{F} = \mathcal{F}$. Il suffit donc de montrer que
$R^qa_{X, *}a_X^{-1}\mathcal{F} = 0$ pour $q > 0$.
Pour cela, nous pouvons utiliser $a_X = \epsilon_X \circ \pi_X$ et la suite
spectrale de Leray (Cohomologie sur les sites, lemme
\ref{sites-cohomology-lemma-relative-Leray}).
D'après le lemme \ref{lemma-comparison-fppf-etale}, on a
$R^i\epsilon_{X, *}(a_X^{-1}\mathcal{F}) = 0$ pour $i > 0$.
On a
$\epsilon_{X, *}a_X^{-1}\mathcal{F} = \pi_X^{-1}\mathcal{F}$
et, d'après le lemme \ref{lemma-cohomological-descent-etale},
$R^j\pi_{X, *}(\pi_X^{-1}\mathcal{F}) = 0$ pour $j > 0$.
Ceci achève la démonstration.
\end{proof}

\begin{lemma}
\label{lemma-compare-cohomology-etale-fppf}
Soient $S$ un schéma et $X$ un espace algébrique sur $S$.
Avec $a_X : \Sh((\textit{Espaces}/X)_{fppf}) \to \Sh(X_\etale)$
défini ci-dessus:
\begin{enumerate}
\item $H^q(X_\etale, \mathcal{F}) = H^q_{fppf}(X, a_X^{-1}\mathcal{F})$
pour un faisceau abélien $\mathcal{F}$ sur $X_\etale$,
\item $H^q(X_\etale, K) = H^q_{fppf}(X, a_X^{-1}K)$ pour $K \in D^+(X_\etale)$.
\end{enumerate}
Exemple: si $A$ est un groupe abélien, alors
$H^q_\etale(X, \underline{A}) = H^q_{fppf}(X, \underline{A})$.
\end{lemma}

\begin{proof}
Ceci découle du lemme \ref{lemma-cohomological-descent-etale-fppf}, d'après
Cohomologie sur les sites, remarque \ref{sites-cohomology-remark-before-Leray}.
\end{proof}

\begin{lemma}
\label{lemma-push-pull-fppf-etale}
Soit $S$ un schéma.
Soit $f : X \to Y$ un morphisme d'espaces algébriques sur $S$.
Il existe alors des diagrammes commutatifs de topos
$$
\xymatrix{
\Sh((\textit{Espaces}/X)_{fppf}) \ar[rr]_{f_{big, fppf}} \ar[d]_{\epsilon_X} & &
\Sh((\textit{Espaces}/Y)_{fppf}) \ar[d]^{\epsilon_Y} \\
\Sh((\textit{Espaces}/X)_\etale) \ar[rr]^{f_{big, \etale}} & &
\Sh((\textit{Espaces}/Y)_\etale)
}
$$
et
$$
\xymatrix{
\Sh((\textit{Espaces}/X)_{fppf}) \ar[rr]_{f_{big, fppf}} \ar[d]_{a_X} & &
\Sh((\textit{Espaces}/Y)_{fppf}) \ar[d]^{a_Y} \\
\Sh(X_\etale) \ar[rr]^{f_{small}} & &
\Sh(Y_\etale)
}
$$
avec $a_X = \pi_X \circ \epsilon_X$ et $a_Y = \pi_X \circ \epsilon_X$.
\end{lemma}

\begin{proof}
Ceci découle immédiatement des définitions des morphismes en jeu; voir
Topologies sur les espaces, section \ref{spaces-topologies-section-fppf},
et la section \ref{section-compare}.
\end{proof}

\begin{lemma}
\label{lemma-proper-push-pull-fppf-etale}
Dans le lemme \ref{lemma-push-pull-fppf-etale}, si $f$ est propre, alors on a
\begin{enumerate}
\item $a_Y^{-1} \circ f_{small, *} = f_{big, fppf, *} \circ a_X^{-1}$, et
\item
$a_Y^{-1}(Rf_{small, *}K) = Rf_{big, fppf, *}(a_X^{-1}K)$
pour $K$ dans $D^+(X_\etale)$ à faisceaux de cohomologie de torsion.
\end{enumerate}
\end{lemma}

\begin{proof}
Démonstration de (1). On pourrait la donner en répétant celle du lemme
\ref{lemma-compare-higher-direct-image-proper}, assertion (1); nous allons
plutôt en déduire le résultat. Comme $\epsilon_{Y, *}$ est le foncteur identité
sur les préfaisceaux sous-jacents, il reflète les isomorphismes. Le lemme
\ref{lemma-comparison-fppf-etale} montre que
$\epsilon_{Y, *} \circ a_Y^{-1} = \pi_Y^{-1}$, et de même pour $X$.
Pour montrer que l'application canonique
$a_Y^{-1}f_{small, *}\mathcal{F} \to f_{big, fppf, *}a_X^{-1}\mathcal{F}$
est un isomorphisme, il suffit de montrer que
\begin{align*}
\pi_Y^{-1}f_{small, *}\mathcal{F}
& =
\epsilon_{Y, *}a_Y^{-1}f_{small, *}\mathcal{F} \\
& \to 
\epsilon_{Y, *}f_{big, fppf, *}a_X^{-1}\mathcal{F} \\
& =
f_{big, \etale, *} \epsilon_{X, *}a_X^{-1}\mathcal{F} \\
& =
f_{big, \etale, *}\pi_X^{-1}\mathcal{F}
\end{align*}
est un isomorphisme. C'est l'assertion (1) du lemme
\ref{lemma-compare-higher-direct-image-proper}.

\medskip\noindent
Pour démontrer (2), nous utilisons
\begin{align*}
R\epsilon_{Y, *}Rf_{big, fppf, *}a_X^{-1}K
& =
Rf_{big, \etale, *}R\epsilon_{X, *}a_X^{-1}K \\
& =
Rf_{big, \etale, *}\pi_X^{-1}K \\
& =
\pi_Y^{-1}Rf_{small, *}K \\
& =
R\epsilon_{Y, *} a_Y^{-1}Rf_{small, *}K
\end{align*}
La première égalité résulte du diagramme commutatif du lemme
\ref{lemma-push-pull-fppf-etale} et de Cohomologie sur les sites, lemme
\ref{sites-cohomology-lemma-derived-pushforward-composition}.
La deuxième est le lemme \ref{lemma-cohomological-descent-etale-fppf}; la
troisième est l'assertion (2) du lemme
\ref{lemma-compare-higher-direct-image-proper}; la quatrième est de nouveau
le lemme \ref{lemma-cohomological-descent-etale-fppf}.
Ainsi le morphisme de changement de base
$a_Y^{-1}(Rf_{small, *}K) \to Rf_{big, fppf, *}(a_X^{-1}K)$
induit un isomorphisme
$$
R\epsilon_{Y, *}a_Y^{-1}Rf_{small, *}K \to
R\epsilon_{Y, *}Rf_{big, fppf, *}a_X^{-1}K
$$
La remarque suivante achève la démonstration: soit
$\alpha : a_Y^{-1}L \to M$, avec $L$ dans $D^+(Y_\etale)$ et $M$ dans
$D^+((\textit{Espaces}/Y)_{fppf})$, telle que $R\epsilon_{Y, *}\alpha$ soit un
isomorphisme. Alors cette application est un isomorphisme. En effet, montrons par
récurrence sur $i$ que $H^i(\alpha)$ est un isomorphisme. Ceci est vrai pour
tout $i$ suffisamment petit. Si cela vaut pour $i \leq i_0$, alors
$R^j\epsilon_{Y, *}H^i(M) = 0$ pour $j > 0$ et $i \leq i_0$ d'après le lemme
\ref{lemma-comparison-fppf-etale}, puisque
$H^i(M) = a_Y^{-1}H^i(L)$ dans cet intervalle. Un argument de suite spectrale
donne alors
$\epsilon_{Y, *}H^{i_0 + 1}(M) = H^{i_0 + 1}(R\epsilon_{Y, *}M)$.
Par conséquent, $\epsilon_{Y, *}H^{i_0 + 1}(M) =
\pi_Y^{-1}H^{i_0 + 1}(L) =
\epsilon_{Y, *}a_Y^{-1}H^{i_0 + 1}(L)$. Il s'ensuit que
$H^{i_0 + 1}(\alpha)$ est un isomorphisme (car $\epsilon_{Y, *}$ reflète les
isomorphismes, puisqu'il est l'identité sur les préfaisceaux sous-jacents),
comme voulu.
\end{proof}

\begin{lemma}
\label{lemma-finite-push-pull-fppf-etale}
Dans le lemme \ref{lemma-push-pull-fppf-etale}, si $f$ est fini, alors
$a_Y^{-1}(Rf_{small, *}K) = Rf_{big, fppf, *}(a_X^{-1}K)$
pour $K$ dans $D^+(X_\etale)$.
\end{lemma}

\begin{proof}
Soit $V \to Y$ un morphisme étale surjectif, où $V$ est un schéma.
Il suffit de montrer que le morphisme de changement de base est un isomorphisme
après restriction à $V$. Nous pouvons donc supposer que $Y$ est un schéma.
Comme le morphisme est fini, donc représentable, nous pouvons alors supposer
que $X$ et $Y$ sont tous deux des schémas. Dans ce cas, le résultat découle du
cas des schémas (Cohomologie étale, lemme
\ref{etale-cohomology-lemma-V-C-all-n-etale-fppf}, assertion (2)), au moyen de
la comparaison des topos exposée à la section \ref{section-api} et, en
particulier, dans le lemme \ref{lemma-compare-cohomology-other-topologies}.
Certains détails sont omis.
\end{proof}

\begin{lemma}
\label{lemma-descent-sheaf-fppf-etale}
Dans le lemme \ref{lemma-push-pull-fppf-etale}, supposons que
$f$ soit plat, localement de présentation finie et surjectif.
Alors le foncteur
$$
\Sh(Y_\etale) \longrightarrow
\left\{
(\mathcal{G}, \mathcal{H}, \alpha)
\middle|
\begin{matrix}
\mathcal{G} \in \Sh(X_\etale),\ \mathcal{H} \in \Sh((\Sch/Y)_{fppf}), \\
\alpha : a_X^{-1}\mathcal{G} \to f_{big, fppf}^{-1}\mathcal{H}
\text{ un isomorphisme}
\end{matrix}
\right\}
$$
qui envoie $\mathcal{F}$ sur
$(f_{small}^{-1}\mathcal{F}, a_Y^{-1}\mathcal{F}, can)$ est une équivalence.
\end{lemma}

\begin{proof}
Le foncteur $a_X^{-1}$ est pleinement fidèle (car
$a_{X, *}a_X^{-1} = \text{id}$ d'après le lemme
\ref{lemma-comparison-fppf-etale}). Le foncteur d'oubli
$(\mathcal{G}, \mathcal{H}, \alpha) \mapsto \mathcal{H}$ identifie donc la
catégorie des triplets à une sous-catégorie pleine de $\Sh((\Sch/Y)_{fppf})$.
De plus, le foncteur $a_Y^{-1}$ est pleinement fidèle; le foncteur du lemme
l'est donc lui aussi.

\medskip\noindent
Supposons donné un recouvrement étale $\{Y_i \to Y\}$.
Soit $f_i : X_i \to Y_i$ le changement de base de $f$.
Notons $f_{ij} = f_i \times f_j : X_i \times_X X_j \to Y_i \times_Y Y_j$.
Affirmation: si le lemme est vrai pour $f_i$ et $f_{ij}$ pour tous $i,j$, il
est vrai pour $f$. En effet, le recouvrement étale donné détermine un
recouvrement étale de l'objet final de chacun des quatre sites
$Y_\etale, X_\etale, (\Sch/Y)_{fppf}, (\Sch/X)_{fppf}$.
Dans chacun des quatre cas, la catégorie des faisceaux est donc équivalente à
la catégorie des données de recollement pour ce recouvrement (Sites, lemme
\ref{sites-lemma-mapping-property-glue}). Un immense diagramme commutatif de
catégories achève alors la démonstration de l'affirmation. Nous omettons les
détails. L'affirmation montre que nous pouvons travailler localement pour la
topologie étale sur $Y$. En particulier, nous pouvons supposer que $Y$ est un
schéma.

\medskip\noindent
Supposons que $Y$ soit un schéma. Choisissons un schéma $X'$ et un morphisme
étale surjectif $s : X' \to X$. Posons $f' = f \circ s : X' \to Y$ et
observons que $f'$ est surjectif, localement de présentation finie et plat.
Affirmation: si le lemme est vrai pour $f'$, il est vrai pour $f$.
En effet, à partir d'un triplet $(\mathcal{G}, \mathcal{H}, \alpha)$ pour $f$,
on obtient par image réciproque par $s$ un triplet
$(s_{small}^{-1}\mathcal{G}, \mathcal{H}, s_{big, fppf}^{-1}\alpha)$ pour $f'$.
Une solution de ce triplet fournit un faisceau $\mathcal{F}$ sur $Y_\etale$
tel que $a_Y^{-1}\mathcal{F} = \mathcal{H}$. D'après le premier paragraphe de
la démonstration, cela signifie que le triplet appartient à l'image essentielle.
Nous sommes ainsi ramenés au cas où $X$ et $Y$ sont tous deux des schémas.
Ce cas découle de Cohomologie étale, lemme
\ref{etale-cohomology-lemma-descent-sheaf-fppf-etale}
au moyen de la discussion de la section \ref{section-api}, et en particulier
du lemme \ref{lemma-compare-cohomology-other-topologies}.
\end{proof}



\section{Comparaison des topologies fppf et étale: modules}
\label{section-fppf-etale-modules}

\noindent
Nous poursuivons l'étude de la section \ref{section-fppf-etale}, en examinant
brièvement ici ce qui se passe pour les faisceaux de modules.

\medskip\noindent
Soient $S$ un schéma et $X$ un espace algébrique sur $S$.
Les morphismes de sites $\epsilon_X$, $\pi_X$ et leur composé $a_X$, introduits
à la section \ref{section-fppf-etale}, se prolongent naturellement en des
morphismes de sites annelés. Le premier s'écrit
$$
\epsilon_X :
((\textit{Espaces}/X)_{fppf}, \mathcal{O})
\longrightarrow
((\textit{Espaces}/X)_\etale, \mathcal{O})
$$
Notons que nous pouvons employer le même symbole pour le faisceau structural,
car ces faisceaux ont effectivement le même préfaisceau sous-jacent. Le second est
$$
\pi_X :
((\textit{Espaces}/X)_\etale, \mathcal{O})
\longrightarrow
(X_\etale, \mathcal{O}_X)
$$
Le troisième est le morphisme
$$
a_X :
((\textit{Espaces}/X)_{fppf}, \mathcal{O})
\longrightarrow
(X_\etale, \mathcal{O}_X)
$$
Rappelons ce que nous savons déjà des modules quasi-cohérents sur ces sites.

\begin{lemma}
\label{lemma-review-quasi-coherent}
Soit $S$ un schéma. Soit $X$ un espace algébrique sur $S$.
Soit $\mathcal{F}$ un $\mathcal{O}_X$-module quasi-cohérent.
\begin{enumerate}
\item La règle
$$
\mathcal{F}^a : (\textit{Espaces}/X)_\etale \longrightarrow \textit{Ab},\quad
(f : Y \to X) \longmapsto \Gamma(Y, f^*\mathcal{F})
$$
satisfait à la condition de faisceau pour les recouvrements fpqc et, a fortiori,
pour les recouvrements fppf et étales,
\item $\mathcal{F}^a = \pi_X^*\mathcal{F}$ sur $(\textit{Espaces}/X)_\etale$,
\item $\mathcal{F}^a = a_X^*\mathcal{F}$ sur $(\textit{Espaces}/X)_{fppf}$,
\item la règle $\mathcal{F} \mapsto \mathcal{F}^a$ définit une équivalence entre
les $\mathcal{O}_X$-modules quasi-cohérents et les modules quasi-cohérents sur
$((\textit{Espaces}/X)_\etale, \mathcal{O})$,
\item la règle $\mathcal{F} \mapsto \mathcal{F}^a$ définit une équivalence entre
les $\mathcal{O}_X$-modules quasi-cohérents et les modules quasi-cohérents sur
$((\textit{Espaces}/X)_{fppf}, \mathcal{O})$,
\item on a $\epsilon_{X, *}a_X^*\mathcal{F} = \pi_X^*\mathcal{F}$
et $a_{X, *}a_X^*\mathcal{F} = \mathcal{F}$,
\item on a $R^i\epsilon_{X, *}(a_X^*\mathcal{F}) = 0$
et $R^ia_{X, *}(a_X^*\mathcal{F}) = 0$ pour $i > 0$.
\end{enumerate}
\end{lemma}

\begin{proof}
L'assertion (1) est une conséquence de la descente fppf des modules
quasi-cohérents. Supposons en effet que $\{f_i : U_i \to U\}$ soit un
recouvrement fpqc dans $(\textit{Espaces}/X)_\etale$. Notons $g : U \to X$ le
morphisme structural. Supposons donnée une famille de sections
$s_i \in \Gamma(U_i , f_i^*g^*\mathcal{F})$ telle que
$s_i|_{U_i \times_U U_j} = s_j|_{U_i \times_U U_j}$. Il faut trouver la section
correspondante $s \in \Gamma(U, g^*\mathcal{F})$. Nous pouvons interpréter
les $s_i$ comme une famille d'applications
$\varphi_i : f_i^*\mathcal{O}_U = \mathcal{O}_{U_i} \to f_i^*g^*\mathcal{F}$
compatibles avec les données de descente canoniques associées aux faisceaux
quasi-cohérents $\mathcal{O}_U$ et $g^*\mathcal{F}$ sur $U$.
Par conséquent, d'après Descente sur les espaces, proposition
\ref{spaces-descent-proposition-fpqc-descent-quasi-coherent}
nous pouvons les descendre (de manière unique) en une application
$\mathcal{O}_U \to g^*\mathcal{F}$, qui fournit notre section $s$.

\medskip\noindent
Nous déduirons les assertions (2) à (7) de l'énoncé correspondant pour les
schémas. Choisissons un recouvrement étale $\{X_i \to X\}_{i \in I}$, où
chaque $X_i$ est un schéma. Observons que $X_i \times_X X_j$ est également un
schéma. Ce recouvrement induit un recouvrement de l'objet final de chacun des
trois sites
$(\textit{Espaces}/X)_{fppf}$, $(\textit{Espaces}/X)_\etale$ et $X_\etale$.
Ainsi, la catégorie des faisceaux sur chacun de ces sites est équivalente à
celle des données de descente pour ces recouvrements; voir Sites, lemme
\ref{sites-lemma-mapping-property-glue}. Les assertions (2) et (3) sont locales
(grâce à l'énoncé de recollement). La quasi-cohérence est une propriété locale;
les assertions (4) et (5) sont donc locales. Les assertions (6) et (7) le sont
manifestement. Il suffit par conséquent de démontrer les assertions (2) à (7)
du lemme lorsque $X$ est un schéma.

\medskip\noindent
Supposons que $X$ soit un schéma. Les plongements
$(\Sch/X)_\etale \subset (\textit{Espaces}/X)_\etale$ et
$(\Sch/X)_{fppf} \subset (\textit{Espaces}/X)_{fppf}$
déterminent des équivalences de topos annelés d'après le lemme
\ref{lemma-compare-cohomology-other-topologies}.
Nous concluons que les assertions (2) à (7) découlent du cas des schémas,
Cohomologie étale, lemme
\ref{etale-cohomology-lemma-review-quasi-coherent}.
Pour transporter la propriété de quasi-cohérence par cette équivalence, on
utilise le fait que la quasi-cohérence est une propriété intrinsèque des
modules, comme expliqué dans Modules sur les sites, section
\ref{sites-modules-section-local}. Quelques détails mineurs sont omis.
\end{proof}

\begin{lemma}
\label{lemma-cohomological-descent-etale-fppf-modules}
Soit $S$ un schéma. Soit $X$ un espace algébrique sur $S$.
Pour $\mathcal{F}$ un $\mathcal{O}_X$-module quasi-cohérent, les applications
$$
\pi_X^*\mathcal{F} \longrightarrow R\epsilon_{X, *}(a_X^*\mathcal{F})
\quad\text{et}\quad
\mathcal{F} \longrightarrow Ra_{X, *}(a_X^*\mathcal{F})
$$
sont des isomorphismes.
\end{lemma}

\begin{proof}
C'est une conséquence immédiate des assertions (6) et (7) du lemme
\ref{lemma-review-quasi-coherent}.
\end{proof}

\begin{lemma}
\label{lemma-vanishing-adequate}
Soit $S$ un schéma. Soit $X$ un espace algébrique sur $S$.
Soit $\mathcal{F}_1 \to \mathcal{F}_2 \to \mathcal{F}_3$
un complexe de $\mathcal{O}_X$-modules quasi-cohérents.
Posons
$$
\mathcal{H}_\etale =
\Ker(\pi_X^*\mathcal{F}_2 \to \pi_X^*\mathcal{F}_3)/
\Im(\pi_X^*\mathcal{F}_1 \to \pi_X^*\mathcal{F}_2)
$$
sur $(\textit{Espaces}/X)_\etale$, et posons
$$
\mathcal{H}_{fppf} =
\Ker(a_X^*\mathcal{F}_2 \to a_X^*\mathcal{F}_3)/
\Im(a_X^*\mathcal{F}_1 \to a_X^*\mathcal{F}_2)
$$
sur $(\textit{Espaces}/X)_{fppf}$.
Alors $\mathcal{H}_\etale = \epsilon_{X, *}\mathcal{H}_{fppf}$ et
$$
H^p_\etale(U, \mathcal{H}_\etale) = H^p_{fppf}(U, \mathcal{H}_{fppf}) = 0
$$
pour $p > 0$ et tout objet affine $U$ de $(\textit{Espaces}/X)_\etale$.
\end{lemma}

\noindent
Il y a mieux: les modules sur $(\textit{Espaces}/X)_{fppf}$ qui, localement pour
la topologie fppf, sont de la forme considérée dans le lemme sont appelés
modules adéquats. Ils forment une sous-catégorie de Serre faible de la catégorie
de tous les $\mathcal{O}$-modules, et leur cohomologie est étudiée dans Modules
adéquats, section \ref{adequate-section-adequate}.

\begin{proof}
Pour tout objet $f : U \to X$ de $(\textit{Espaces}/X)_\etale$, considérons la
restriction $\mathcal{H}_\etale|_{U_\etale}$ de $\mathcal{H}_\etale$ à
$U_\etale$ par le foncteur $i_f^* = i_f^{-1}$ étudié à la section
\ref{section-compare}. Le faisceau $\mathcal{H}_\etale|_{U_\etale}$ est égal à
l'homologie du complexe $f^*\mathcal{F}_\bullet$ en degré $1$.
Ceci résulte de l'égalité $i_f \circ \pi_X = f$ comme morphismes de sites
annelés $U_\etale \to X_\etale$. En particulier,
$\mathcal{H}_\etale|_{U_\etale}$ est un $\mathcal{O}_U$-module quasi-cohérent.
Soit ensuite $g : V \to U$ un morphisme plat dans
$(\textit{Espaces}/X)_\etale$. Comme
$$
i_{f \circ g}^* \circ \pi_X^* = (f \circ g)^* = g^* \circ f^*
$$
comme morphismes de sites $V_\etale \to X_\etale$, et comme $g$ est plat, donc
$g^*$ exact, nous obtenons
$$
\mathcal{H}_\etale|_{V_\etale} =
g^*\left(\mathcal{H}_\etale|_{U_\etale}\right)
$$
Ces préparatifs achevés, nous pouvons démontrer le lemme.

\medskip\noindent
Soit $\mathcal{U} = \{g_i : U_i \to U\}_{i \in I}$ un recouvrement fppf, avec
$f : U \to X$ comme ci-dessus. La propriété de faisceau vaut pour
$\mathcal{H}_\etale$ et le recouvrement $\mathcal{U}$, d'après l'assertion (1)
du lemme \ref{lemma-review-quasi-coherent}, appliquée à
$\mathcal{H}_\etale|_{U_\etale}$, et d'après ce qui précède.
Ainsi $\mathcal{H}_\etale$ est déjà un faisceau fppf, ce qui signifie que
$\mathcal{H}_{fppf}$ est égal à $\mathcal{H}_\etale$ comme préfaisceau.
En particulier,
$\mathcal{H}_\etale = \epsilon_{X, *}\mathcal{H}_{fppf}$.

\medskip\noindent
Enfin, pour démontrer l'annulation, nous utilisons Cohomologie sur les sites,
lemme \ref{sites-cohomology-lemma-cech-vanish-collection}. Soit $\mathcal{B}$
la classe des objets affines de $(\textit{Espaces}/X)_{fppf}$ et soit
$\text{Cov}$ l'ensemble des recouvrements fppf finis
$\mathcal{U} = \{U_i \to U\}_{i = 1, \ldots, n}$ où $U$ et les $U_i$ sont affines.
On a
$$
{\check H}^p(\mathcal{U}, \mathcal{H}_\etale) =
{\check H}^p(\mathcal{U}, \left(\mathcal{H}_\etale|_{U_\etale}\right)^a)
$$
car, d'après le premier paragraphe, les valeurs de $\mathcal{H}_\etale$ sur les
schémas affines $U_{i_0} \times_U \ldots \times_U U_{i_p}$, plats sur $U$,
coïncident avec les valeurs de l'image réciproque du module quasi-cohérent
$\mathcal{H}_\etale|_{U_\etale}$. L'annulation résulte alors de Descente, lemme
\ref{descent-lemma-standard-covering-Cech-quasi-coherent}.
Ceci achève la démonstration.
\end{proof}

\begin{lemma}
\label{lemma-cohomological-descent-etale-fppf-modules-unbounded}
Soit $S$ un schéma. Soit $X$ un espace algébrique sur $S$.
Pour $K \in D_\QCoh(\mathcal{O}_X)$, les applications
$$
L\pi_X^*K \longrightarrow R\epsilon_{X, *}(La_X^*K)
\quad\text{et}\quad
K \longrightarrow Ra_{X, *}(La_X^*K)
$$
sont des isomorphismes. Ici,
$a_X : \Sh((\textit{Espaces}/X)_{fppf}) \to \Sh(X_\etale)$ est défini ci-dessus.
\end{lemma}

\begin{proof}
La question est locale pour la topologie étale sur $X$; nous pouvons donc
supposer que $X$ est affine, disons $X = \Spec(A)$. Alors
$D_\QCoh(\mathcal{O}_X) = D(A)$ d'après Catégories dérivées des espaces, lemme
\ref{spaces-perfect-lemma-derived-quasi-coherent-small-etale-site}
et Catégories dérivées des schémas, lemme
\ref{perfect-lemma-affine-compare-bounded}.
Nous pouvons donc choisir un complexe K-plat de $A$-modules $K^\bullet$ dont
le complexe correspondant $\mathcal{K}^\bullet$ de $\mathcal{O}_X$-modules
quasi-cohérents représente $K$.
Nous affirmons que $\mathcal{K}^\bullet$ est un complexe K-plat de
$\mathcal{O}_X$-modules.

\medskip\noindent
Démonstration de l'affirmation. D'après Catégories dérivées des schémas, lemme
\ref{perfect-lemma-affine-K-flat}
le complexe $\widetilde{K}^\bullet$ est K-plat sur le schéma
$(\Spec(A), \mathcal{O}_{\Spec(A)})$.
Ensuite, notons que $\mathcal{K}^\bullet = \epsilon^*\widetilde{K}^\bullet$,
où $\epsilon$ est défini comme dans Catégories dérivées des espaces, lemme
\ref{spaces-perfect-lemma-derived-quasi-coherent-small-etale-site}
; le complexe $\mathcal{K}^\bullet$ est donc K-plat d'après Cohomologie sur les
sites, lemme \ref{sites-cohomology-lemma-pullback-K-flat-points}, et le fait que
le site étale d'un schéma possède suffisamment de points (Cohomologie étale,
remarques \ref{etale-cohomology-remarks-enough-points}).

\medskip\noindent
D'après l'affirmation, on a
$La_X^*K = a_X^*\mathcal{K}^\bullet$ et
$L\pi_X^*K = \pi_X^*\mathcal{K}^\bullet$.
Comme la première partie de la démonstration montre que l'image réciproque
$a_X^*\mathcal{K}^n$ du module quasi-cohérent est acyclique pour
$\epsilon_{X, *}$, resp.\ $a_{X, *}$, le lemme d'acyclicité de Leray ne
permet-il pas aussitôt de conclure? En fait, non, car ce lemme ne s'applique
qu'aux complexes bornés inférieurement. Nous montrerons toutefois au paragraphe
suivant que le résultat découle bien du cas borné inférieurement, parce que
notre complexe est la limite dérivée de complexes bornés inférieurement de
modules quasi-cohérents.

\medskip\noindent
Les faisceaux de cohomologie de $\pi_X^*\mathcal{K}^\bullet$ et de
$a_X^*\mathcal{K}^\bullet$ ont des groupes de cohomologie supérieure nuls sur
les objets affines de $(\textit{Espaces}/X)_\etale$, d'après le lemme
\ref{lemma-vanishing-adequate}. Par conséquent, on a
$$
L\pi_X^*K = R\lim \tau_{\geq -n}(L\pi_X^*K)
\quad\text{et}\quad
La_X^*K = R\lim \tau_{\geq -n}(La_X^*K)
$$
d'après Cohomologie sur les sites, lemme
\ref{sites-cohomology-lemma-is-limit-dimension}.

\medskip\noindent
Démonstration de $L\pi_X^*K = R\epsilon_{X, *}(La_X^*\mathcal{F})$.
D'après ce qui précède, on a
$$
R\epsilon_{X, *}La_X^*K =
R\lim R\epsilon_{X, *}(\tau_{\geq -n}(La_X^*K))
$$
d'après Cohomologie sur les sites, lemme
\ref{sites-cohomology-lemma-Rf-commutes-with-Rlim}.
Notons que $\tau_{\geq -n}(La_X^*K)$ est représenté par
$\tau_{\geq -n}(a_X^*\mathcal{K}^\bullet)$, qui peut différer de
$a_X^*(\tau_{\geq -n}\mathcal{K}^\bullet)$.
Mais les systèmes
$$
\{\tau_{\geq -n}(a_X^*\mathcal{K}^\bullet)\}_{n \geq 1}
\quad\text{et}\quad
\{a_X^*(\tau_{\geq -n}\mathcal{K}^\bullet)\}_{n \geq 1}
$$
sont manifestement isomorphes comme pro-systèmes.
D'après le lemme d'acyclicité de Leray (Catégories dérivées, lemme
\ref{derived-lemma-leray-acyclicity}) et la première partie du lemme, on a
$$
R\epsilon_{X, *}(a_X^*(\tau_{\geq -n}\mathcal{K}^\bullet)) =
\pi_X^*(\tau_{\geq -n}\mathcal{K}^\bullet)
$$
Nous pouvons ensuite utiliser le fait que les systèmes
$$
\{\tau_{\geq -n}(\pi_X^*\mathcal{K}^\bullet)\}_{n \geq 1}
\quad\text{et}\quad
\{\pi_X^*(\tau_{\geq -n}\mathcal{K}^\bullet)\}_{n \geq 1}
$$
sont isomorphes comme pro-systèmes. Enfin, regroupons toutes les égalités:
\begin{align*}
R\epsilon_{X, *}La_X^*K
& =
R\epsilon_{X, *} (R\lim \tau_{\geq -n}(La_X^*K)) \\
& =
R\lim R\epsilon_{X, *}(\tau_{\geq -n}(La_X^*K)) \\
& =
R\lim R\epsilon_{X, *}(\tau_{\geq -n}(a_X^*\mathcal{K}^\bullet)) \\
& =
R\lim R\epsilon_{X, *}(a_X^*(\tau_{\geq -n}\mathcal{K}^\bullet)) \\
& =
R\lim \pi_X^*(\tau_{\geq -n}\mathcal{K}^\bullet) \\
& =
R\lim \tau_{\geq -n}(\pi_X^*\mathcal{K}^\bullet) \\
& =
R\lim \tau_{\geq -n}(L\pi_X^*K) \\
& =
L\pi_X^*K
\end{align*}
Dans les quatrième et sixième égalités, nous avons utilisé que des pro-systèmes
isomorphes ont le même $R\lim$ (un petit détail est omis).
On peut éviter cette étape en utilisant les résultats plus précis sur la
cohomologie des termes du complexe $\tau_{\geq -n}a_X^*\mathcal{K}^\bullet$
démontrés au lemme \ref{lemma-vanishing-adequate}; ils donnent directement
$R\epsilon_{X, *}(\tau_{\geq -n}(a_X^*\mathcal{K}^\bullet)) =
\tau_{\geq -n}(\pi_X^*\mathcal{K}^\bullet)$.

\medskip\noindent
L'égalité $K = Ra_{X, *}(La_X^*\mathcal{F})$ se démontre exactement de la même
manière, en utilisant à la dernière étape l'égalité
$K = R\lim \tau_{\geq -n}K$, donnée par Catégories dérivées des espaces,
lemme \ref{spaces-perfect-lemma-nice-K-injective}.
\end{proof}








\section{Comparaison des topologies ph et étale}
\label{section-ph-etale}

\noindent
Cette section est l'analogue de Cohomologie étale, section
\ref{etale-cohomology-section-ph-etale}.

\medskip\noindent
Soient $S$ un schéma et $X$ un espace algébrique sur $S$.
Sur la catégorie $\textit{Espaces}/X$, considérons les topologies ph et étale.
Le foncteur identité
$(\textit{Espaces}/X)_\etale \to (\textit{Espaces}/X)_{ph}$
est continu, puisque tout recouvrement étale est un recouvrement ph d'après
Topologies sur les espaces, lemme
\ref{spaces-topologies-lemma-zariski-etale-smooth-syntomic-fppf-ph}.
Il définit donc un morphisme de sites
$$
\epsilon_X :
(\textit{Espaces}/X)_{ph} \longrightarrow (\textit{Espaces}/X)_\etale
$$
par application de Sites, proposition \ref{sites-proposition-get-morphism}.
Notons que $\epsilon_{X, *}$ est le foncteur identité sur les préfaisceaux
sous-jacents et que $\epsilon_X^{-1}$ associe à un faisceau étale son faisceau
associé pour la topologie ph.
Considérons le morphisme de sites
$$
\pi_X : (\textit{Espaces}/X)_\etale \longrightarrow X_{spaces, \etale}
$$
qui compare les grands et petits sites étales; voir la section \ref{section-compare}.
Le composé détermine un morphisme de sites
$$
a_X = \pi_X \circ \epsilon_X :
(\textit{Espaces}/X)_{ph}
\longrightarrow
X_{spaces, \etale}
$$
Si $\mathcal{H}$ est un faisceau abélien sur $(\textit{Espaces}/X)_{ph}$,
nous noterons $H^n_{ph}(U, \mathcal{H})$ la cohomologie de
$\mathcal{H}$ sur un objet $U$ de $(\textit{Espaces}/X)_{ph}$.

\begin{lemma}
\label{lemma-comparison-ph-etale}
Soit $S$ un schéma. Soit $X$ un espace algébrique sur $S$.
\begin{enumerate}
\item Pour $\mathcal{F} \in \Sh(X_\etale)$, on a
$\epsilon_{X, *}a_X^{-1}\mathcal{F} = \pi_X^{-1}\mathcal{F}$
et $a_{X, *}a_X^{-1}\mathcal{F} = \mathcal{F}$.
\item Pour $\mathcal{F} \in \textit{Ab}(X_\etale)$ de torsion, on a
$R^i\epsilon_{X, *}(a_X^{-1}\mathcal{F}) = 0$ pour $i > 0$.
\end{enumerate}
\end{lemma}

\begin{proof}
On a $a_X^{-1}\mathcal{F} = \epsilon_X^{-1} \pi_X^{-1}\mathcal{F}$.
D'après le lemme \ref{lemma-describe-pullback}, le faisceau étale
$\pi_X^{-1}\mathcal{F}$ est un faisceau pour la topologie ph; il est donc égal
à $a_X^{-1}\mathcal{F}$ (puisque l'image réciproque par $\epsilon_X$ est donnée
par le faisceau associé pour la topologie ph). Rappelons en outre que
$\epsilon_{X, *}$ est le foncteur identité sur les préfaisceaux sous-jacents.
L'assertion (1) résulte alors immédiatement de la description explicite de
$\pi_X^{-1}$ donnée au lemme \ref{lemma-describe-pullback}.

\medskip\noindent
Nous démontrerons (2) en nous ramenant au cas des schémas; voir l'assertion (1)
de Cohomologie étale, lemme
\ref{etale-cohomology-lemma-V-C-all-n-etale-ph}.
Cette méthode « fonctionne clairement », puisque tout espace algébrique est
localement un schéma pour la topologie étale. Les détails sont donnés ci-dessous,
mais nous invitons vivement le lecteur à omettre la démonstration.

\medskip\noindent
Pour un faisceau abélien $\mathcal{H}$ sur $(\textit{Espaces}/X)_{ph}$,
l'image directe supérieure $R^p\epsilon_{X, *}\mathcal{H}$ est le faisceau
associé au préfaisceau $U \mapsto H^p_{ph}(U, \mathcal{H})$
sur $(\textit{Espaces}/X)_\etale$. Voir Cohomologie sur les sites, lemme
\ref{sites-cohomology-lemma-higher-direct-images}.
Comme tout objet de $(\textit{Espaces}/X)_\etale$ admet un recouvrement par des
schémas, il suffit de démontrer que, pour tout schéma $U/X$ et tout
$\xi \in H^p_{ph}(U, a_X^{-1}\mathcal{F})$, il existe un recouvrement étale
$\{U_i \to U\}$ tel que la restriction de $\xi$ à chaque $U_i$ soit nulle.
On a
\begin{align*}
H^p_{ph}(U, a_X^{-1}\mathcal{F})
& =
H^p((\textit{Espaces}/U)_{ph}, (a_X^{-1}\mathcal{F})|_{\textit{Espaces}/U}) \\
& =
H^p((\Sch/U)_{ph}, (a_X^{-1}\mathcal{F})|_{\Sch/U})
\end{align*}
où la seconde identification est le lemme
\ref{lemma-compare-cohomology-other-topologies}, et la première un fait général
sur la restriction (Cohomologie sur les sites, lemme
\ref{sites-cohomology-lemma-cohomology-of-open}).
En comparant le premier paragraphe et le résultat correspondant dans le cas
des schémas (Cohomologie étale, lemme
\ref{etale-cohomology-lemma-describe-pullback-pi-ph})
nous concluons que le faisceau $(a_X^{-1}\mathcal{F})|_{\Sch/U}$
coïncide avec l'image réciproque par la « version pour les schémas de $a_U$ ».
Nous pouvons donc trouver un recouvrement étale $\{U_i \to U\}$ tel que notre
classe s'annule dans
$H^p((\Sch/U_i)_{ph}, (a_X^{-1}\mathcal{F})|_{\Sch/U_i})$
pour tout $i$; voir Cohomologie étale, lemme
\ref{etale-cohomology-lemma-V-C-all-n-etale-ph}
(l'énoncé précis à utiliser ici est que $V_n$ vaut pour tout $n$, ce qui est
l'assertion (2) dans le cas des schémas).
En revenant en arrière (au moyen des mêmes formules que ci-dessus, appliquées
cette fois à $U_i$), nous concluons que la restriction de $\xi$ à $U_i$
est nulle, comme souhaité.
\end{proof}

\noindent
Le travail difficile effectué dans le cas des schémas montre maintenant que
les cohomologies étale et ph coïncident pour les faisceaux abéliens de torsion
provenant du petit site étale.

\begin{lemma}
\label{lemma-cohomological-descent-etale-ph}
Soit $S$ un schéma. Soit $X$ un espace algébrique sur $S$.
Pour $K \in D^+(X_\etale)$ à faisceaux de cohomologie de torsion, les applications
$$
\pi_X^{-1}K \longrightarrow R\epsilon_{X, *}a_X^{-1}K
\quad\text{et}\quad
K \longrightarrow Ra_{X, *}a_X^{-1}K
$$
sont des isomorphismes, où
$a_X : \Sh((\textit{Espaces}/X)_{ph}) \to \Sh(X_\etale)$ est défini ci-dessus.
\end{lemma}

\begin{proof}
Nous démontrons seulement le second énoncé; le premier est plus facile et se
démontre exactement de la même manière. On se ramène au cas où $K$ est donné
par un seul faisceau abélien de torsion. En effet, représentons $K$ par un
complexe borné inférieurement $\mathcal{F}^\bullet$ de faisceaux abéliens de
torsion. C'est possible d'après Cohomologie sur les sites, lemme
\ref{sites-cohomology-lemma-torsion}.
Le cas d'un faisceau donne
$\mathcal{F}^n = a_{X, *} a_X^{-1} \mathcal{F}^n$
et les faisceaux $R^qa_{X, *}a_X^{-1}\mathcal{F}^n$ sont nuls pour $q > 0$.
Le lemme d'acyclicité de Leray (Catégories dérivées, lemme
\ref{derived-lemma-leray-acyclicity}), appliqué à
$a_X^{-1}\mathcal{F}^\bullet$ et au foncteur $a_{X, *}$, permet de conclure.
Supposons désormais $K = \mathcal{F}$, où $\mathcal{F}$ est un faisceau
abélien de torsion.

\medskip\noindent
Le lemme \ref{lemma-comparison-ph-etale} donne
$a_{X, *}a_X^{-1}\mathcal{F} = \mathcal{F}$. Il suffit donc de montrer que
$R^qa_{X, *}a_X^{-1}\mathcal{F} = 0$ pour $q > 0$.
Pour cela, nous pouvons utiliser $a_X = \epsilon_X \circ \pi_X$ et la suite
spectrale de Leray (Cohomologie sur les sites, lemme
\ref{sites-cohomology-lemma-relative-Leray}).
D'après le lemme \ref{lemma-comparison-ph-etale}, on a
$R^i\epsilon_{X, *}(a_X^{-1}\mathcal{F}) = 0$ pour $i > 0$.
On a
$\epsilon_{X, *}a_X^{-1}\mathcal{F} = \pi_X^{-1}\mathcal{F}$
et, d'après le lemme \ref{lemma-cohomological-descent-etale},
$R^j\pi_{X, *}(\pi_X^{-1}\mathcal{F}) = 0$ pour $j > 0$.
Ceci achève la démonstration.
\end{proof}

\begin{lemma}
\label{lemma-compare-cohomology-etale-ph}
Soient $S$ un schéma et $X$ un espace algébrique sur $S$.
Avec $a_X : \Sh((\textit{Espaces}/X)_{ph}) \to \Sh(X_\etale)$
défini ci-dessus:
\begin{enumerate}
\item $H^q(X_\etale, \mathcal{F}) = H^q_{ph}(X, a_X^{-1}\mathcal{F})$
pour un faisceau abélien de torsion $\mathcal{F}$ sur $X_\etale$,
\item $H^q(X_\etale, K) = H^q_{ph}(X, a_X^{-1}K)$ pour $K \in D^+(X_\etale)$
à faisceaux de cohomologie de torsion.
\end{enumerate}
Exemple: si $A$ est un groupe abélien de torsion, alors
$H^q_\etale(X, \underline{A}) = H^q_{ph}(X, \underline{A})$.
\end{lemma}

\begin{proof}
Ceci découle du lemme \ref{lemma-cohomological-descent-etale-ph}, d'après
Cohomologie sur les sites, remarque \ref{sites-cohomology-remark-before-Leray}.
\end{proof}

\begin{lemma}
\label{lemma-push-pull-ph-etale}
Soit $S$ un schéma.
Soit $f : X \to Y$ un morphisme d'espaces algébriques sur $S$.
Il existe alors des diagrammes commutatifs de topos
$$
\xymatrix{
\Sh((\textit{Espaces}/X)_{ph}) \ar[rr]_{f_{big, ph}} \ar[d]_{\epsilon_X} & &
\Sh((\textit{Espaces}/Y)_{ph}) \ar[d]^{\epsilon_Y} \\
\Sh((\textit{Espaces}/X)_\etale) \ar[rr]^{f_{big, \etale}} & &
\Sh((\textit{Espaces}/Y)_\etale)
}
$$
et
$$
\xymatrix{
\Sh((\textit{Espaces}/X)_{ph}) \ar[rr]_{f_{big, ph}} \ar[d]_{a_X} & &
\Sh((\textit{Espaces}/Y)_{ph}) \ar[d]^{a_Y} \\
\Sh(X_\etale) \ar[rr]^{f_{small}} & &
\Sh(Y_\etale)
}
$$
avec $a_X = \pi_X \circ \epsilon_X$ et $a_Y = \pi_X \circ \epsilon_X$.
\end{lemma}

\begin{proof}
Ceci découle immédiatement des définitions des morphismes en jeu; voir
Topologies sur les espaces, section \ref{spaces-topologies-section-ph},
et la section \ref{section-compare}.
\end{proof}

\begin{lemma}
\label{lemma-proper-push-pull-ph-etale}
Dans le lemme \ref{lemma-push-pull-ph-etale}, si $f$ est propre, alors on a
\begin{enumerate}
\item $a_Y^{-1} \circ f_{small, *} = f_{big, ph, *} \circ a_X^{-1}$, et
\item
$a_Y^{-1}(Rf_{small, *}K) = Rf_{big, ph, *}(a_X^{-1}K)$
pour $K$ dans $D^+(X_\etale)$ à faisceaux de cohomologie de torsion.
\end{enumerate}
\end{lemma}

\begin{proof}
Démonstration de (1). On pourrait la donner en répétant celle du lemme
\ref{lemma-compare-higher-direct-image-proper}, assertion (1); nous allons
plutôt en déduire le résultat. Comme $\epsilon_{Y, *}$ est le foncteur identité
sur les préfaisceaux sous-jacents, il reflète les isomorphismes. Le lemme
\ref{lemma-comparison-ph-etale} montre que
$\epsilon_{Y, *} \circ a_Y^{-1} = \pi_Y^{-1}$, et de même pour $X$.
Pour montrer que l'application canonique
$a_Y^{-1}f_{small, *}\mathcal{F} \to f_{big, ph, *}a_X^{-1}\mathcal{F}$
est un isomorphisme, il suffit de montrer que
\begin{align*}
\pi_Y^{-1}f_{small, *}\mathcal{F}
& =
\epsilon_{Y, *}a_Y^{-1}f_{small, *}\mathcal{F} \\
& \to 
\epsilon_{Y, *}f_{big, ph, *}a_X^{-1}\mathcal{F} \\
& =
f_{big, \etale, *} \epsilon_{X, *}a_X^{-1}\mathcal{F} \\
& =
f_{big, \etale, *}\pi_X^{-1}\mathcal{F}
\end{align*}
est un isomorphisme. C'est l'assertion (1) du lemme
\ref{lemma-compare-higher-direct-image-proper}.

\medskip\noindent
Pour démontrer (2), nous utilisons
\begin{align*}
R\epsilon_{Y, *}Rf_{big, ph, *}a_X^{-1}K
& =
Rf_{big, \etale, *}R\epsilon_{X, *}a_X^{-1}K \\
& =
Rf_{big, \etale, *}\pi_X^{-1}K \\
& =
\pi_Y^{-1}Rf_{small, *}K \\
& =
R\epsilon_{Y, *} a_Y^{-1}Rf_{small, *}K
\end{align*}
La première égalité résulte du diagramme commutatif du lemme
\ref{lemma-push-pull-ph-etale} et de Cohomologie sur les sites, lemme
\ref{sites-cohomology-lemma-derived-pushforward-composition}.
La deuxième est le lemme \ref{lemma-cohomological-descent-etale-ph}; la
troisième est l'assertion (2) du lemme
\ref{lemma-compare-higher-direct-image-proper}; la quatrième est de nouveau
le lemme \ref{lemma-cohomological-descent-etale-ph}.
Ainsi le morphisme de changement de base
$a_Y^{-1}(Rf_{small, *}K) \to Rf_{big, ph, *}(a_X^{-1}K)$
induit un isomorphisme
$$
R\epsilon_{Y, *}a_Y^{-1}Rf_{small, *}K \to
R\epsilon_{Y, *}Rf_{big, ph, *}a_X^{-1}K
$$
La remarque suivante achève la démonstration: soit
$\alpha : a_Y^{-1}L \to M$, avec $L$ dans $D^+(Y_\etale)$ à faisceaux de
cohomologie de torsion et $M$ dans $D^+((\textit{Espaces}/Y)_{ph})$.
Si $R\epsilon_{Y, *}\alpha$ est un isomorphisme, alors $\alpha$ est un
isomorphisme. En effet, montrons par récurrence sur $i$ que $H^i(\alpha)$ est un
isomorphisme. Ceci est vrai pour tout $i$ suffisamment petit. Si cela vaut pour
$i \leq i_0$, alors $R^j\epsilon_{Y, *}H^i(M) = 0$ pour $j > 0$ et
$i \leq i_0$ d'après le lemme \ref{lemma-comparison-ph-etale}, puisque
$H^i(M) = a_Y^{-1}H^i(L)$ dans cet intervalle. Un argument de suite spectrale
donne alors
$\epsilon_{Y, *}H^{i_0 + 1}(M) = H^{i_0 + 1}(R\epsilon_{Y, *}M)$.
Par conséquent, $\epsilon_{Y, *}H^{i_0 + 1}(M) =
\pi_Y^{-1}H^{i_0 + 1}(L) =
\epsilon_{Y, *}a_Y^{-1}H^{i_0 + 1}(L)$. Il s'ensuit que
$H^{i_0 + 1}(\alpha)$ est un isomorphisme (car $\epsilon_{Y, *}$ reflète les
isomorphismes, puisqu'il est l'identité sur les préfaisceaux sous-jacents),
comme voulu.
\end{proof}













\begin{multicols}{2}[\section{Autres chapitres}]
\noindent
Préliminaires
\begin{enumerate}
\item \hyperref[introduction-section-phantom]{Introduction}
\item \hyperref[conventions-section-phantom]{Conventions}
\item \hyperref[sets-section-phantom]{Théorie des ensembles}
\item \hyperref[categories-section-phantom]{Catégories}
\item \hyperref[topology-section-phantom]{Topologie}
\item \hyperref[sheaves-section-phantom]{Faisceaux sur les espaces}
\item \hyperref[sites-section-phantom]{Sites et faisceaux}
\item \hyperref[stacks-section-phantom]{Champs}
\item \hyperref[fields-section-phantom]{Corps}
\item \hyperref[algebra-section-phantom]{Algèbre commutative}
\item \hyperref[brauer-section-phantom]{Groupes de Brauer}
\item \hyperref[homology-section-phantom]{Algèbre homologique}
\item \hyperref[derived-section-phantom]{Catégories dérivées}
\item \hyperref[simplicial-section-phantom]{Méthodes simpliciales}
\item \hyperref[more-algebra-section-phantom]{Compléments d'algèbre}
\item \hyperref[smoothing-section-phantom]{Lissage des morphismes d'anneaux}
\item \hyperref[modules-section-phantom]{Faisceaux de modules}
\item \hyperref[sites-modules-section-phantom]{Modules sur les sites}
\item \hyperref[injectives-section-phantom]{Objets injectifs}
\item \hyperref[cohomology-section-phantom]{Cohomologie des faisceaux}
\item \hyperref[sites-cohomology-section-phantom]{Cohomologie sur les sites}
\item \hyperref[dga-section-phantom]{Algèbre différentielle graduée}
\item \hyperref[dpa-section-phantom]{Algèbre à puissances divisées}
\item \hyperref[sdga-section-phantom]{Faisceaux différentiels gradués}
\item \hyperref[hypercovering-section-phantom]{Hyperrecouvrements}
\end{enumerate}
Schémas
\begin{enumerate}
\setcounter{enumi}{25}
\item \hyperref[schemes-section-phantom]{Schémas}
\item \hyperref[constructions-section-phantom]{Constructions de schémas}
\item \hyperref[properties-section-phantom]{Propriétés des schémas}
\item \hyperref[morphisms-section-phantom]{Morphismes de schémas}
\item \hyperref[coherent-section-phantom]{Cohomologie des schémas}
\item \hyperref[divisors-section-phantom]{Diviseurs}
\item \hyperref[limits-section-phantom]{Limites de schémas}
\item \hyperref[varieties-section-phantom]{Variétés}
\item \hyperref[topologies-section-phantom]{Topologies sur les schémas}
\item \hyperref[descent-section-phantom]{Descente}
\item \hyperref[perfect-section-phantom]{Catégories dérivées de schémas}
\item \hyperref[more-morphisms-section-phantom]{Compléments sur les morphismes}
\item \hyperref[flat-section-phantom]{Compléments sur la platitude}
\item \hyperref[groupoids-section-phantom]{Schémas en groupoïdes}
\item \hyperref[more-groupoids-section-phantom]{Compléments sur les schémas en groupoïdes}
\item \hyperref[etale-section-phantom]{Morphismes étales de schémas}
\end{enumerate}
Thèmes de théorie des schémas
\begin{enumerate}
\setcounter{enumi}{41}
\item \hyperref[chow-section-phantom]{Homologie de Chow}
\item \hyperref[intersection-section-phantom]{Théorie de l'intersection}
\item \hyperref[pic-section-phantom]{Schémas de Picard des courbes}
\item \hyperref[weil-section-phantom]{Théories cohomologiques de Weil}
\item \hyperref[adequate-section-phantom]{Modules adéquats}
\item \hyperref[dualizing-section-phantom]{Complexes dualisants}
\item \hyperref[duality-section-phantom]{Dualité pour les schémas}
\item \hyperref[discriminant-section-phantom]{Discriminants et différentes}
\item \hyperref[derham-section-phantom]{Cohomologie de de Rham}
\item \hyperref[local-cohomology-section-phantom]{Cohomologie locale}
\item \hyperref[algebraization-section-phantom]{Géométrie algébrique et formelle}
\item \hyperref[curves-section-phantom]{Courbes algébriques}
\item \hyperref[resolve-section-phantom]{Résolution des surfaces}
\item \hyperref[models-section-phantom]{Réduction semi-stable}
\item \hyperref[functors-section-phantom]{Foncteurs et morphismes}
\item \hyperref[equiv-section-phantom]{Catégories dérivées de variétés}
\item \hyperref[pione-section-phantom]{Groupes fondamentaux de schémas}
\item \hyperref[etale-cohomology-section-phantom]{Cohomologie étale}
\item \hyperref[crystalline-section-phantom]{Cohomologie cristalline}
\item \hyperref[proetale-section-phantom]{Cohomologie pro-étale}
\item \hyperref[relative-cycles-section-phantom]{Cycles relatifs}
\item \hyperref[more-etale-section-phantom]{Compléments de cohomologie étale}
\item \hyperref[trace-section-phantom]{La formule des traces}
\end{enumerate}
Espaces algébriques
\begin{enumerate}
\setcounter{enumi}{64}
\item \hyperref[spaces-section-phantom]{Espaces algébriques}
\item \hyperref[spaces-properties-section-phantom]{Propriétés des espaces algébriques}
\item \hyperref[spaces-morphisms-section-phantom]{Morphismes d'espaces algébriques}
\item \hyperref[decent-spaces-section-phantom]{Espaces algébriques décents}
\item \hyperref[spaces-cohomology-section-phantom]{Cohomologie des espaces algébriques}
\item \hyperref[spaces-limits-section-phantom]{Limites d'espaces algébriques}
\item \hyperref[spaces-divisors-section-phantom]{Diviseurs sur les espaces algébriques}
\item \hyperref[spaces-over-fields-section-phantom]{Espaces algébriques sur des corps}
\item \hyperref[spaces-topologies-section-phantom]{Topologies sur les espaces algébriques}
\item \hyperref[spaces-descent-section-phantom]{Descente et espaces algébriques}
\item \hyperref[spaces-perfect-section-phantom]{Catégories dérivées d'espaces}
\item \hyperref[spaces-more-morphisms-section-phantom]{Compléments sur les morphismes d'espaces}
\item \hyperref[spaces-flat-section-phantom]{Platitude sur les espaces algébriques}
\item \hyperref[spaces-groupoids-section-phantom]{Groupoïdes dans les espaces algébriques}
\item \hyperref[spaces-more-groupoids-section-phantom]{Compléments sur les groupoïdes dans les espaces}
\item \hyperref[bootstrap-section-phantom]{Amorçage}
\item \hyperref[spaces-pushouts-section-phantom]{Sommes amalgamées d'espaces algébriques}
\end{enumerate}
Thèmes de géométrie
\begin{enumerate}
\setcounter{enumi}{81}
\item \hyperref[spaces-chow-section-phantom]{Groupes de Chow des espaces}
\item \hyperref[groupoids-quotients-section-phantom]{Quotients de groupoïdes}
\item \hyperref[spaces-more-cohomology-section-phantom]{Compléments sur la cohomologie des espaces}
\item \hyperref[spaces-simplicial-section-phantom]{Espaces simpliciaux}
\item \hyperref[spaces-duality-section-phantom]{Dualité pour les espaces}
\item \hyperref[formal-spaces-section-phantom]{Espaces algébriques formels}
\item \hyperref[restricted-section-phantom]{Algébrisation des espaces formels}
\item \hyperref[spaces-resolve-section-phantom]{Résolution des surfaces revisitée}
\end{enumerate}
Théorie des déformations
\begin{enumerate}
\setcounter{enumi}{89}
\item \hyperref[formal-defos-section-phantom]{Théorie formelle des déformations}
\item \hyperref[defos-section-phantom]{Théorie des déformations}
\item \hyperref[cotangent-section-phantom]{Le complexe cotangent}
\item \hyperref[examples-defos-section-phantom]{Problèmes de déformation}
\end{enumerate}
Champs algébriques
\begin{enumerate}
\setcounter{enumi}{93}
\item \hyperref[algebraic-section-phantom]{Champs algébriques}
\item \hyperref[examples-stacks-section-phantom]{Exemples de champs}
\item \hyperref[stacks-sheaves-section-phantom]{Faisceaux sur les champs algébriques}
\item \hyperref[criteria-section-phantom]{Critères de représentabilité}
\item \hyperref[artin-section-phantom]{Axiomes d'Artin}
\item \hyperref[quot-section-phantom]{Espaces de Quot et de Hilbert}
\item \hyperref[stacks-properties-section-phantom]{Propriétés des champs algébriques}
\item \hyperref[stacks-morphisms-section-phantom]{Morphismes de champs algébriques}
\item \hyperref[stacks-limits-section-phantom]{Limites de champs algébriques}
\item \hyperref[stacks-cohomology-section-phantom]{Cohomologie des champs algébriques}
\item \hyperref[stacks-perfect-section-phantom]{Catégories dérivées de champs}
\item \hyperref[stacks-introduction-section-phantom]{Introduction aux champs algébriques}
\item \hyperref[stacks-more-morphisms-section-phantom]{Compléments sur les morphismes de champs}
\item \hyperref[stacks-geometry-section-phantom]{La géométrie des champs}
\end{enumerate}
Thèmes de théorie des espaces de modules
\begin{enumerate}
\setcounter{enumi}{107}
\item \hyperref[moduli-section-phantom]{Champs de modules}
\item \hyperref[moduli-curves-section-phantom]{Espaces de modules de courbes}
\end{enumerate}
Divers
\begin{enumerate}
\setcounter{enumi}{109}
\item \hyperref[examples-section-phantom]{Exemples}
\item \hyperref[exercises-section-phantom]{Exercices}
\item \hyperref[guide-section-phantom]{Guide bibliographique}
\item \hyperref[desirables-section-phantom]{Desiderata}
\item \hyperref[coding-section-phantom]{Style de programmation}
\item \hyperref[obsolete-section-phantom]{Obsolète}
\item \hyperref[fdl-section-phantom]{Licence de documentation libre GNU}
\item \hyperref[index-section-phantom]{Index généré automatiquement}
\end{enumerate}
\end{multicols}


\bibliography{my}
\bibliographystyle{amsalpha}

\end{document}
